\documentclass[11pt,a4paper]{article}
\usepackage[T1]{fontenc}
\usepackage[utf8]{inputenc}
\usepackage{lmodern,microtype}
\usepackage[margin=26mm,headheight=15pt]{geometry}
\usepackage{amsmath,amssymb,amsthm,mathtools}
\usepackage{booktabs,longtable,array,enumitem,xcolor,listings,fancyhdr,xurl}
\usepackage[colorlinks=true,linkcolor=blue!45!black,citecolor=blue!45!black,urlcolor=blue!50!black]{hyperref}
\definecolor{ink}{HTML}{17324D}
\definecolor{paper}{HTML}{F5F7FA}
\hypersetup{pdftitle={Inherited foundations for concise iPRA},pdfauthor={ProveIt research sources; compilation prepared for Vladimir Reshetnikov},pdfsubject={Attributed scalar, logical, and progressive-closure foundation proofs}}
\pagestyle{fancy}\fancyhf{}
\fancyhead[L]{\small\textsc{iPRA: inherited foundations}}
\fancyhead[R]{\small November 11, 2026}
\fancyfoot[C]{\thepage}
\setlength{\parindent}{0pt}\setlength{\parskip}{4pt}
\setlength{\emergencystretch}{3em}
\setlist{itemsep=3pt,topsep=5pt}
\newtheorem{theorem}{Theorem}[section]
\newtheorem{lemma}[theorem]{Lemma}
\newtheorem{proposition}[theorem]{Proposition}
\newtheorem{corollary}[theorem]{Corollary}
\theoremstyle{definition}\newtheorem{definition}[theorem]{Definition}
\newtheorem{remark}[theorem]{Remark}
\newcommand{\ipra}{\mathrm{iPRA}}
\newcommand{\Teq}{T_{\mathrm{eq}}}
\newcommand{\eqv}{\dashv\vdash}
\newcommand{\FV}{\operatorname{FV}}
\newcommand{\llbracket}{\mathopen{[\![}}
\newcommand{\rrbracket}{\mathclose{]\!]}}
\lstset{basicstyle=\ttfamily\footnotesize,columns=fullflexible,keepspaces=true,breaklines=true,showstringspaces=false,frame=single,framerule=.3pt,rulecolor=\color{black!20},backgroundcolor=\color{paper},aboveskip=9pt,belowskip=9pt}

\newcommand{\IPRA}{\mathrm{iPRA}}
\newcommand{\Q}{\mathsf Q}
\newcommand{\D}{\mathsf D}
\newcommand{\code}[1]{\texttt{\detokenize{#1}}}
\newcommand{\file}[1]{{\small\path{#1}}}
\title{\color{ink}\textbf{Inherited foundations for concise iPRA}\\[6pt]\large Attributed companion source compilation}
\author{ProveIt research sources; compilation prepared for Vladimir Reshetnikov}
\date{November 11, 2026}
\begin{document}
\maketitle
This companion supplies the existing mathematical results used by the main
report. The new target-directed rule and lexical validation theorem are in
\emph{Target-directed closure and smaller typed checking for iPRA}; they are
not claimed as results of the historical sections reproduced here.

The source snapshot is ProveIt commit
\texttt{76697262071ada91f29c14d9cf3801a64dacb994}, directory
\path{Logic/PrimitiveRecursiveArithmetic/ConciseIPRA}. The five scalar
foundation modules have been checked byte-for-byte against their Git blob
identities in \path{Research/certified_zero/reference/}; the receipt is
\path{evidence/foundation_identity.json}. The zero-source theorem was
retrieved directly from the pinned \texttt{source\_dispatch} package. The
progressive closure sections are reproduced from the incoming
\texttt{ipra\_uniform\_progress\_2026-11-11.zip}, whose Git blob is
\texttt{2b8392d6036daf0272b8697b549557aefc9e1cad}. The upstream MIT No
Attribution license is retained in \path{reference/LICENSE-ProveIt.txt}.

Historical phrases such as ``new'', ``current'', and ``the new calculus''
in reproduced sections refer to their original reports. In particular,
$\mathcal R$ in the progressive part is the inherited orbit calculus,
not the target-directed $\mathcal T$ of the main report. Its forward
translation is called $T$ in the historical text and $F$ in the main
report. The native conditional theory is a different theory with the
uniform reading $Q B=B(0)\to\forall y B(y)$.

This is a compilation of written proofs, not a fresh verification of the
upstream Lean project. The full-language progressive roundtrip theorem
is an inherited written theorem; the native/scalar/iPRA bridge also has
upstream Lean endpoint files, which were inspected but not rebuilt here.

\setcounter{tocdepth}{1}
{\small\tableofcontents}
\clearpage
\part{Conditional native logic and scalar arithmetic}

% BEGIN reference/foundations/inherited_choice_basis.tex
% New mathematical result: complementary introduction choices.
% Requires amsmath/amssymb/amsthm and theorem, lemma, proposition environments.
% The inherited arithmetic-to-iPRA bridge must be supplied separately.

\section{Complementary introductions with one local assumption}
\label{sec:choice-introduction}

The old paired introduction retains both the zero instance and an
induction instance when its matrix is an equation. Both together are
unnecessary. A certificate may choose either one, provided the induction
choice is restricted to a literal equation. This changes neither the
formula language nor the meaning of a formula. The omitted combination
is recovered by a short nested introduction--elimination construction.

\subsection{Complete native syntax and rule inventory}

The scalar numerical language is
\[
 t::=x\mid0\mid St\mid I(n,a;u.h).
\]
Only $h$ binds the state variable $u$; $n$ and $a$ do not.
The update may contain fixed numerical parameters. The computation
axioms in every retained context are
\[
 I(0,a;u.h)=a,\qquad
 I(Sn,a;u.h)=h[I(n,a;u.h)/u].
 \tag{I0, IS}\label{eq:choice-computations}
\]
The argument also applies unchanged to the full fixed-description PR
language with its projection, composition, and recursion equations.
No numerical variable ranges over function descriptions, and no uniform
evaluator is postulated.

Native formulas and reusable context clauses are
\[
 A::=(s=t)\mid\mathsf D_{x,y}A,\qquad
 \langle A;y.B(y)\rangle,\quad y\notin\operatorname{FV}(A).
\]
Bound variables are renamed before capture-free substitution.
The two variables $x,y$ in a paired formula are distinct.
An ordinary assumption $F$ is the constant clause
$\langle H;v.F\rangle$, where $v$ is absent from $F$ and $H$ is the
fixed computation theorem $I(0,0;u.0)=0$. Thus ordinary hypothesis
use is a clause Use preceded by this computation proof.

There are computation axioms and exactly the following inference
patterns:
\[
\frac{\Gamma\vdash A}{\Gamma\vdash B(t)}
\quad\mathsf{Use}
\quad(\langle A;y.B(y)\rangle\in\Gamma),
\]
\[
\frac{\Gamma\vdash\mathsf D_{x,y}B(x,y)\qquad
      \Gamma,\langle B(a,0);y.B(a,y)\rangle\vdash F}
     {\Gamma\vdash F}
\quad\mathsf{DE},
\qquad
\frac{\Gamma\vdash s=t\qquad\Gamma\vdash F(s)}
     {\Gamma\vdash F(t)}
\quad\mathsf{Tr},
\]
\[
\frac{\Gamma,B(t,\eta_\varepsilon(z))\vdash B(t,Sz)}
     {\Gamma\vdash\mathsf D_{x,y}B(x,y)}
\quad\mathsf{DI}_{\varepsilon},
\qquad
\eta_0(z)=0,\quad\eta_1(z)=z.
\tag{Choice DI}\label{eq:choice-di}
\]
Here $\varepsilon$ is a certificate choice in $\{0,1\}$, and
$\varepsilon=1$ is permitted only when $B$ has a literal equation
root. An equation beneath a binder is not a literal equation.
Call $\varepsilon=0$ the \emph{case choice} and $\varepsilon=1$
the \emph{progress choice}. Both retain exactly one ordinary
formula assumption. Compound matrices have only the case choice.

In DI the eigenparameter $z$ is fresh for $\Gamma$, the witness
$t$, and every fixed parameter of $B$. In DE the eigenparameter
$a$ is fresh for $\Gamma$, the conclusion $F$, and every fixed
parameter of $B$. In transport, $F(v)$ is any fixed native
formula template. Displayed local assumptions and clauses are
discharged. There are no primitive reflexivity, coverage,
separation, ordinary logical, or additional induction rules.

We call this calculus $\mathcal C$. The original comparison
calculus $\mathcal N$ has identical syntax, computation,
Use, DE, and transport, but its introduction is
\[
\frac{\Gamma,B(t,0),[B(t,z)]_{\mathrm{at}}\vdash B(t,Sz)}
     {\Gamma\vdash\mathsf D_{x,y}B(x,y)}.
\tag{Old DI}\label{eq:choice-old-di}
\]
The bracket supplies $B(t,z)$ exactly at a literal equation.

\subsection{Exact equivalence and the unused-witness construction}

\begin{lemma}[Structural admissibility]
\label{lem:choice-structural}
Weakening, exchange, contraction, renaming, ordinary-assumption cut,
and capture-free numerical substitution are admissible in $\mathcal C$.
\end{lemma}
\begin{proof}
Induct on derivations. Rename local eigenparameters before inserting
proofs or terms. Numerical substitution preserves a formula's root
constructor, so an allowed progress choice remains allowed.
The selected choice does not change under substitution or cut.
For ordinary-assumption cut, replace each use of its constant clause
by the supplied proof. Computation schemas are closed under
capture-free numerical substitution. These facts verify every case.
\end{proof}

\begin{theorem}[Identity-language equivalence]
\label{thm:choice-equivalence}
For every retained clause context $\Gamma$ and native formula $F$,
\[
 \Gamma\vdash_{\mathcal C}F
 \quad\Longleftrightarrow\quad
 \Gamma\vdash_{\mathcal N}F.
\]
The transformations are effective on finite derivations and leave
all numerical and formula syntax unchanged.
\end{theorem}
\begin{proof}
Each new introduction is an instance of the old one after weakening
by the omitted local assumption. For the case choice at an equation,
add the induction instance. For the progress choice, add the zero
instance. Compound case introduction is already the old rule.
Thus every new derivation translates to an old one.

For the reverse direction, transform an old derivation recursively.
Only its atomic DI case needs a construction. Suppose its matrix is
the equation $E(x,y)$ and its transformed premise is
\[
 \pi:\quad\Gamma,E(t,0),E(t,z)\vdash E(t,Sz).
\]
The outer new inference will be case DI with witness $t$. It is
therefore enough to prove $E(t,Sz)$ under
$\Gamma,E(t,0)$.

Rename the eigenvariable in $\pi$ to a fresh $v$. Choose a new
variable $u$ absent from $E(t,y)$ and all retained data. Under
$\Gamma,E(t,0)$, progress DI with witness zero introduces
\[
             \mathsf D_{u,y}E(t,y)
\]
from exactly the renamed premise $\pi[v/z]$.
The bound coordinate $u$ is deliberately unused.

Apply DE to this paired formula. Because $u$ is unused, the opened
clause is exactly
\[
               \langle E(t,0);y.E(t,y)\rangle;
\]
it contains no new existential witness. The retained ordinary
assumption $E(t,0)$ supplies its antecedent. Use at $Sz$ gives
$E(t,Sz)$, and DE closes. Finally the outer case DI introduces
the original $\mathsf D_{x,y}E(x,y)$ with witness $t$.
The original premise proof $\pi$ has been inserted once.

An old compound introduction already is the new case choice.
All other rule cases are unchanged. This completes the induction
on the original derivation.
\end{proof}

The unused inner coordinate is essential. Introducing
$\mathsf D_{x,y}E(x,y)$ internally and then eliminating it would
give an arbitrary fresh witness, not necessarily the original
$t$. The construction instead keeps $t$ as a fixed parameter
of $E(t,y)$ and binds only the unused coordinate $u$.

\begin{proposition}[Linear bound on inference nodes]
\label{prop:choice-node-bound}
Count each computation axiom, Use, DI, DE, and transport as one
inference node, expanding ordinary-assumption use as a computation
of $H$ followed by Use. If an old derivation has $N$ nodes and
$r$ atomic DI nodes, the displayed transformation has exactly
\[
                   N+5r\leq 6N
\]
nodes.
\end{proposition}
\begin{proof}
Each atomic old DI node is replaced by six nodes: outer case DI,
inner progress DI, DE, Use of the opened clause, Use of the
retained base assumption, and the computation of $H$ needed for
that last ordinary-assumption use. Its old premise subtree
occurs once. All other inference nodes are copied once.
\end{proof}

This counts inference-tree nodes. It is not a bound on serialized
JSON bytes, formula sizes, checker runtime, or the number of
consumed packet cells when hypothesis closure is implicit.
In the formula-context Q/D presentation, the corresponding
replacement has outer case DI, inner progress DI, DE, QE,
and two Hyp nodes, giving the same six-node count.

\subsection{Induction and numerical cases separate cleanly}

\begin{lemma}[Contextual equation induction]
\label{lem:choice-ei}
The progress choice, DE, and Use derive
\[
\frac{\Gamma\vdash E(0)\qquad
      \Gamma,E(z)\vdash E(Sz)}
     {\Gamma\vdash E(t)}
\tag{EI}\label{eq:choice-ei}
\]
for every equation $E$, with the usual freshness condition.
\end{lemma}
\begin{proof}
Choose an unused witness coordinate $u$. Progress DI introduces
$\mathsf D_{u,y}E(y)$ from the given step. DE opens the clause
$\langle E(0);y.E(y)\rangle$. Supply its antecedent using the given
base proof and use the clause at $t$. Close DE.
\end{proof}

\begin{lemma}[Cases, equality, and binder congruence]
\label{lem:choice-preinduction}
The case choice derives, for arbitrary native formulas,
\[
\frac{\Gamma\vdash F(0)\qquad\Gamma\vdash F(Sz)}
     {\Gamma\vdash F(t)}
\tag{NC}\label{eq:choice-nc}
\]
with the usual freshness condition. Equality reflexivity,
symmetry, transitivity, numerical congruence, and uniform
matrix interderivability under $\mathsf D$ are derivable
without using the progress choice.
\end{lemma}
\begin{proof}
For cases, use case DI with an unused witness coordinate and
matrix $F(y)$, ignoring its retained base in the successor
premise. DE opens $\langle F(0);y.F(y)\rangle$; the given zero
proof and Use at $t$ finish.

For reflexivity put $r=I(0,t;u.0)$. Use the computation equation
$r=t$ as both transport premises with motive $v=t$ to obtain
$t=t$. The usual explicit equality motives give symmetry,
transitivity, and numerical congruence.

For binder congruence, open the old paired formula at a fresh
witness $a$ and introduce the new one by case DI at $a$.
Convert its retained base to the old base, use the opened
old clause at the successor, and convert back. Ordinary cut
closes the construction, and the reverse conversion gives
the converse. No induction is used.
\end{proof}

The primitive progress choice is still restricted to equations.
Allowing it at arbitrary compound matrices would, by the
argument for \eqref{eq:choice-ei}, supply induction for
arbitrary formulas and change the intended theory.

\subsection{The entire ordinary logical interface}
\label{subsec:choice-logic}

Here is a predecessor-free reconstruction of the ordinary logical
operations. It makes explicit that the case choice supplies the
logical interface independently of progress induction.

Let $C(v,a,b):=I(v,a;u.b)$ with $u$ unused. Computation gives
$C(0,a,b)=a$ and $C(Sv,a,b)=b$. Define the gate on fixed finite
formula syntax by
\[
 G_v(s=t):=(C(v,0,s)=C(v,0,t)),\qquad
 G_v(\mathsf D_{x,y}B):=\mathsf D_{x,y}G_vB,
\]
renaming bound variables away from the control. Equality,
binder congruence, and external structural induction give
\[
 \vdash G_0A,\qquad G_{Sv}A\dashv\vdash A,\qquad
 A\dashv\vdash\mathsf D_{u,y}G_yA
 \quad(u,y\notin\operatorname{FV}(A)).
\tag{Gate and padding}\label{eq:choice-gate}
\]
For padding, case DI proves the forward direction. DE, the
zero-gate theorem, and Use at $1$ prove the reverse direction.

Define a selector between two arbitrary formulas by
\[
 M_v(s=t,u=w):=(C(v,s,u)=C(v,t,w))
\]
and, with aligned fresh paired variables, by
\[
\begin{aligned}
 M_v(\mathsf D_{x,y}A,\mathsf D_{x,y}B)
   &:=\mathsf D_{x,y}M_v(A,B),\\
 M_v(E,\mathsf D_{x,y}B)
   &:=\mathsf D_{x,y}M_v(G_yE,B),\\
 M_v(\mathsf D_{x,y}A,E)
   &:=\mathsf D_{x,y}M_v(A,G_yE),
\end{aligned}
\]
where $E$ is an equation. Every call decreases maximum input
binder depth. Padding and congruence prove
$M_0(A,B)\dashv\vdash A$ and
$M_{Sv}(A,B)\dashv\vdash B$.

Put $\delta(v)=C(v,1,0)$ and $p(z)=I(z,0;u.\delta(u))$.
Only the fixed computations $p(1)=1$ and $p(2)=0$ are needed.
The full logical macro inventory is
\[
\begin{aligned}
 A\Rightarrow_*B&:=\mathsf D_{u,z}M_z(A,B),\\
 A\wedge_*B&:=\mathsf D_{u,z}G_zM_{p(z)}(A,B),\\
 \exists_*x\,A(x)&:=\mathsf D_{x,y}G_yA(x),\\
 A\vee_*B&:=\exists_*x\,M_x(A,B),\\
 \forall_*x\,A(x)&:=
        A(0)\wedge_*\mathsf D_{u,z}G_zA(z),\\
 \bot_*&:=(0=1).
\end{aligned}
\tag{Logical macros}\label{eq:choice-macros}
\]
Newly bound variables are fresh; $u$ is unused.

Implication introduction obtains its assumption from the zero
selector and supplies the consequent at successors, then uses
case DI. Elimination opens its clause and uses it at $1$.
Existential introduction uses its witness and the successor
gate; elimination opens at a fresh witness, supplies the
zero-gate theorem, and uses the clause at $1$.

Disjunction introduction uses witnesses $0$ and $1$.
For elimination, discharge the two branch proofs as
$A\Rightarrow_*F$ and $B\Rightarrow_*F$. NC applied to
$M_z(A,B)\Rightarrow_*F$ gives this implication at every
control: the two endpoint cases follow from the branch
implications. Existential elimination provides $M_a(A,B)$
and the implication gives $F$.

Given $A$ and $B$, selector endpoints and NC give $M_v(A,B)$
at an arbitrary control. Substitute $v=p(Sz)$ and use the
successor gate for conjunction introduction. For projections,
open its zero-gated clause and specialize at $1$ and $2$;
the two fixed computations give $B$ and $A$, respectively.

These rules prove existential numerical coverage:
\[
     (x=0)\vee_*\exists_*v\,(x=Sv).
\tag{Coverage}\label{eq:choice-coverage}
\]
Use NC. Zero takes the left injection of reflexivity, and a
successor takes the right injection with the evident witness.
Coverage gives equality cases even when the retained context
contains $x$; its elimination adds $x=0$ or $x=Sv$ rather
than silently changing that retained context.

Universal introduction follows from a fresh-parameter proof
by substituting zero and successors and forming the two
conjuncts. For elimination, take its $A(0)$ conjunct and
open the paired conjunct. Its base is a zero-gate theorem.
Coverage of the target chooses the zero instance from
$A(0)$ or a successor instance from the opened clause
and gate removal.

Finally, under $0=1$, computation and transport give
$s=C(0,s,t)=C(1,s,t)=t$ for arbitrary numerical terms.
Every equation follows, and external induction on formula
syntax with case DI gives every native formula. This
supplies falsity elimination. No excluded middle for
quantified formulas has entered the construction.

\subsection{Whole-domain realizations and both formula composites}

For a standard first-order reading, let $T_{\rm eq}$ have
intuitionistic equality logic, the same numerical signature
and computation equations, numerical coverage, separation
$(0=1)\to\bot$, and contextual equation induction.
Its interpretation of the native syntax is
\[
 (s=t)^\circ=(s=t),\qquad
 (\mathsf D_{x,y}B)^\circ
   =\exists x\bigl(B^\circ(x,0)\to\forall y B^\circ(x,y)\bigr),
\]
\[
 \langle A;y.B\rangle^\circ
       =A^\circ\to\forall y B^\circ.
\]
The domain is the entire numerical domain and interpreted
equality is actual equality.

Case DI is sound in this realization by numerical coverage.
Progress DI is sound by contextual equation induction under
the conditional base. Use, DE, transport, and computation
have their displayed ordinary meanings. Conversely, the
native macros derive every ordinary logical rule,
\eqref{eq:choice-ei} derives equation induction, and
\eqref{eq:choice-coverage} derives coverage. Separation
translates to the identity implication
$(0=1)\Rightarrow_*(0=1)$.

For the essential paired formula composite,
\[
 \mathsf D_{x,y}B
 \dashv\vdash
 \exists_*x\bigl(B(x,0)\Rightarrow_*\forall_*y\,B(x,y)\bigr).
\tag{Paired composite}\label{eq:choice-paired-composite}
\]
Forwards, DE opens its clause at a fresh $a$; under $B(a,0)$
that clause supplies $B(a,y)$ at a fresh arbitrary $y$.
Introduce the universal, implication, and existential
macros and close DE. Backwards, eliminate the existential
at a fresh $a$. Introduce the paired formula with witness
$a$ by case DI: the retained base, supplied implication,
and universal elimination at $Sz$ give its successor premise.

The ordinary interpretations of the six logical macros
have their ordinary round trips. For conjunction, its true
base and the instances $1,2$ recover both operands;
the converse uses selector cases. The universal has all
successor instances in its paired conjunct and zero in
its first conjunct; coverage gives exactly the ordinary
universal. Implication, disjunction, and existential
quantification use their selector or gate endpoints.
Falsity uses the explicit ordinary separation bridge.
External induction on finite formulas, together with
\eqref{eq:choice-paired-composite}, gives both formula
composites. Clauses correspond in both directions to
$A\Rightarrow_*\forall_*yB$ by the derived logical rules.

The numerical composites are literal identities with
internally provable graph $x=y$. No standard-numeral
restriction, retract, quotient, or equality reflection
assumption is being inserted at this stage. The separately
supplied arithmetic bridge identifies $T_{\rm eq}$ with
full iPRA; the present change preserves that bridge
unchanged because Theorem~\ref{thm:choice-equivalence}
preserves every native sequent.

\subsection{What the theorem does and does not claim}

The new inference interface retains one ordinary assumption
at every paired introduction. It preserves the entire
consequence relation, not just equations or closed
theorems. Consequently the established full-iPRA
interpretations and equiconsistency transfer unchanged.

This result is different from retaining only progress
introduction and strengthening closure elimination.
Here the original elimination rule is retained, and
the case and progress introductions coexist. It is also
different from allowing arbitrary compound progress,
which would supply stronger induction.

A complete-checker comparison must still count the
certificate choice, its equation-root restriction,
substitution, context handling, computation, and the
public input/output interface. The mathematical theorem
alone is not a byte-count or minimality theorem.
The inference-node bound is distinct from source-size
measurements and from serialized certificate size.
The proof is a written universal proof; finite regression
tests do not constitute a Lean/Rocq formalization of it.

% END reference/foundations/inherited_choice_basis.tex


% BEGIN reference/foundations/scalar_predecessor.tex
% Adapted with attribution from the pinned unparenthesized one-hole archive.
\section{Internal scalar predecessor bootstrap}
\label{sec:hybrid-bootstrap}

This scalar bootstrap is inherited from the supplied scalar development~\cite{erasure} and is reproduced to keep the scalar checker self-contained. It is not a new contribution of this paper. The proof uses a factor map fixing odd states and a saturated-iteration construction of predecessor.

The available arithmetic terms are $0$, $S$, and
\[
 I(n,a;u.h),\qquad
 I(0,a;u.h)=a,\qquad
 I(Sn,a;u.h)=h[I(n,a;u.h)/u].
\]
The update binds only $u$.  All other variables in it are fixed numerical
parameters.  We use equality reasoning, the already reconstructed
zero/successor cases, and induction on a single literal equation.
Every induction below keeps its other parameters fixed.  Predecessor,
successor injectivity, pairing, and induction on arbitrary
quantifier-free formulas are not assumptions of the argument.

\subsection{Elementary terms and two finite cycles}

Write $1=S0$ and use larger numerals as successor abbreviations.  Define
\[
 C(n,a,b)=I(n,a;u.b),\qquad
 Z(n)=C(n,1,0),\qquad
 a+b=I(b,a;u.Su).
\]
The update variable of $C$ is absent from its body.  Consequently
\[
 C(0,a,b)=a,\qquad C(Sn,a,b)=b.
\]
Numerical cases show that $Z(n)$ is $0$ or $1$.
The usual commutative-monoid equations for addition follow by literal
equation induction.  For example, $0+b=b$ and $Sa+b=S(a+b)$ are
proved by induction on $b$; these yield commutativity by a further
equation induction, and associativity follows by induction on the last
summand.  We use these equations to reorder finite sums.  The notation
$2a$ means $a+a$ and $3a$ means $a+a+a$; no multiplication operation
is needed for this bootstrap.

Put
\[
 N(n)=I(n,0;u.Z(u)),\qquad
 j(u)=Z(u)+2N(u),\qquad
 M(n)=I(n,0;u.j(u)).
\]
The range of $N$ is contained in $\{0,1\}$: its zero value is $0$,
and its successor value is a value of $Z$.  This uses numerical cases
on its counter, not induction on a disjunction.  Hence
\begin{equation}
 N(n)+Z(N(n))=1,\qquad
 Z(Z(N(n)))=N(n),\qquad N(n+2)=N(n).
 \label{eq:hybrid-parity}
\end{equation}
Here the last equation follows from two successor computations and
the preceding equation.  Literal-equation inductions give
\begin{equation}
 N(2n)=0,\qquad N(2n+1)=1,\qquad N(3n)=N(n).
 \label{eq:hybrid-parity-scaling}
\end{equation}
For the first equation the successor step adds two to the argument of
$N$; the second follows by one computation.  In the third, increasing
$n$ adds three to the left argument and therefore applies $Z$ once to
both sides of the induction equation.

Finite computations now give
\[
 j(0)=1,\qquad j(1)=2,\qquad j(2)=0.
\]
The proof below needs just these three values and
$M(Sn)=j(M(n))$; it does not assume a division algorithm or a theorem
characterizing $M$ as remainder modulo three.

\subsection{A ramp with a global scaling identity}

Define
\[
 G(n)=I(n,0;u.u+M(u)+1).
\]
The central invariant is
\begin{equation}
 M(G(n))=N(n).
 \label{eq:hybrid-ramp-invariant}
\end{equation}
Its zero case is immediate.  For its successor step, retain the
single induction equation and split the Boolean value $N(n)$.
If it is $0$, the next $G$ state is $G(n)+1$; its $M$ value is
$j(0)=1$, equal to $N(Sn)$.  If it is $1$, the next $G$ state
is $G(n)+2$; two computations of $M$ change its value from $1$
to $2$ and then to $0$, again equal to $N(Sn)$.  This proves
the invariant by literal-equation induction.

It follows that
\[
 G(Sn)=G(n)+N(n)+1,\qquad G(n+2)=G(n)+3.
\]
Equation induction therefore yields
\begin{equation}
 G(2n)=3n.
 \label{eq:hybrid-ramp-double}
\end{equation}
A separate literal-equation induction proves
\begin{equation}
 G(3n)=3G(n)+N(n).
 \label{eq:hybrid-ramp-triple}
\end{equation}
For completeness, its successor calculation is
\begin{align*}
 G(3Sn)
 &=G(3n)+N(3n)+4
   =3G(n)+2N(n)+4,\\
 3G(Sn)+N(Sn)
 &=3G(n)+3N(n)+3+Z(N(n)).
\end{align*}
The expressions agree in both Boolean cases for $N(n)$.
The base is $0=0$.  In particular,
\eqref{eq:hybrid-ramp-triple} was proved internally; it was not
inferred from an intended floor-function interpretation of $G$.

Set
\[
 H(u)=C(N(u),G(u),u).
\]
The crucial choice is to retain $u$ when its parity is $1$.
The preceding equations give
\begin{equation}
 H(2u)=3u,\qquad H(1)=1,\qquad H(3u)=3H(u).
 \label{eq:hybrid-factor-map}
\end{equation}
To check the last identity, use $N(3u)=N(u)$ and split the Boolean
$N(u)$.  At $0$, both sides are $3G(u)$ by
\eqref{eq:hybrid-ramp-triple}; at $1$, both are $3u$.
This argument does not require that an arbitrary number has already
been represented as $2v$ or $2v+1$.

\subsection{A comparator from scalar iteration}

For every fixed update $h$, equation induction on $n$, keeping $m,a$
and all parameters fixed, proves the semigroup identity
\[
 I(m+n,a;h)=I(n,I(m,a;h);h).
\]
The base is the zero computation and the step applies $h$ to the
induction equation. This identity will also be used in the later
pairing construction.


For a fixed update $h$, equation induction on $y$, with its initializer
$a$ and all parameters of $h$ fixed, proves the rotation identity
\begin{equation}
 I(Sy,a;h)=I(y,h(a);h).
 \label{eq:hybrid-rotation}
\end{equation}
Both sides have base $h(a)$, and their successor equations apply the
same $h$ to the previous equality.  Here and below $I(y,a;h)$ is
an abbreviation for one fixed finite scalar term, not quantification
over functions.

By \eqref{eq:hybrid-factor-map}, another equation induction gives
\begin{equation}
 I(y,3a;H)=3I(y,a;H).
 \label{eq:hybrid-iterator-scaling}
\end{equation}
Its base is $3a=3a$.  In the step the fixed-parameter induction
equation rewrites the new left side as
\[
 H(3I(y,a;H))=3H(I(y,a;H)).
\]
The initializer parameter stays $a$ throughout this induction.
Equation induction also gives $I(y,1;H)=1$, since $H(1)=1$.

Define
\[
 E(x)=I(x,1;u.2u),\qquad
 L(x,y)=N\bigl(I(y,E(x);H)\bigr).
\]
Then the following entire recursion specification is provable:
\begin{equation}
 L(0,y)=1,\qquad L(Sx,0)=0,\qquad L(Sx,Sy)=L(x,y).
 \label{eq:hybrid-comparator}
\end{equation}
The first equation uses $E(0)=1$ and the fixed-point identity for
$H$.  The second uses $E(Sx)=2E(x)$ and $N(2a)=0$.
For the third,
\begin{align*}
 L(Sx,Sy)
 &=N\bigl(I(Sy,2E(x);H)\bigr)\\
 &=N\bigl(I(y,H(2E(x));H)\bigr)\\
 &=N\bigl(I(y,3E(x);H)\bigr)\\
 &=N\bigl(3I(y,E(x);H)\bigr)\\
 &=N\bigl(I(y,E(x);H)\bigr)=L(x,y).
\end{align*}
The successive steps use rotation, $H(2a)=3a$, iterator scaling,
and $N(3a)=N(a)$.  Thus $L$ has the comparison equations at all
arguments internally.  Its standard-model meaning as the indicator
of $x\le y$ is motivation, not a premise of the following argument.

\subsection{Saturated iteration gives predecessor}

We now need only the three equations
\eqref{eq:hybrid-comparator}.  Define
\[
 B_m(u)=C(L(Su,m),u,Su),\qquad
 W(t,m)=I(t,0;u.B_m(u)).
\]
The subscript in $B_m$ records a fixed parameter; $B_m$ is an open
term abbreviation, not a function variable.

Numerical cases on $e$ prove the selector equation
\[
 C(e,Sa,Sb)=S C(e,a,b).
\]
Equivalently one may use equation induction on $e$, whose induction
hypothesis is unused.  Therefore
\begin{align}
 B_{Sm}(Su)
 &=C(L(SSu,Sm),Su,SSu)\notag\\
 &=C(L(Su,m),Su,SSu)=S B_m(u),
 \label{eq:hybrid-conjugacy}\\
 B_{Sm}(0)&=C(L(S0,Sm),0,S0)=C(L(0,m),0,S0)=S0.
 \label{eq:hybrid-conjugacy-base}
\end{align}

Literal-equation induction on $t$, with $m$ fixed, proves
\begin{equation}
 W(St,Sm)=S W(t,m).
 \label{eq:hybrid-saturated-shift}
\end{equation}
The base is \eqref{eq:hybrid-conjugacy-base}.  The successor step is
\begin{align*}
 W(SSt,Sm)
 &=B_{Sm}(W(St,Sm))\\
 &=B_{Sm}(S W(t,m))\\
 &=S B_m(W(t,m))=S W(St,m).
\end{align*}
Only the second line uses the current induction hypothesis, at the
same fixed $m$.

Two subsequent equation inductions give
\begin{equation}
 L(Sn,n)=0,\qquad W(n,n)=n.
 \label{eq:hybrid-diagonals}
\end{equation}
The first uses the comparator recurrence.  The second uses
\eqref{eq:hybrid-saturated-shift} in its successor step.  The
substitution $t=m=n$ is made into that already proved universally
parameterized equation; no parameter of an active induction hypothesis
is changed.  Consequently
\begin{equation}
 W(Sn,n)=B_n(n)=C(L(Sn,n),n,Sn)=n.
 \label{eq:hybrid-extra-step}
\end{equation}

Finally define the finite scalar term
\begin{equation}
 P(m)=I\bigl(m,0;u.C(L(SSu,m),u,Su)\bigr).
 \label{eq:hybrid-predecessor-term}
\end{equation}
Its zero equation is immediate.  To justify its successor equation
without assuming a separate extensionality rule for iterator bodies,
put
\[
 V(t,m)=I\bigl(t,0;u.C(L(SSu,Sm),u,Su)\bigr).
\]
The comparator recurrence identifies this update pointwise with
$B_m(u)$.  Equation induction on $t$, with $m$ fixed, gives
$V(t,m)=W(t,m)$: their base values are $0$, and the successor step
uses the pointwise update equation and then the induction equation.
Now \eqref{eq:hybrid-extra-step} yields
\[
 \boxed{\qquad P(0)=0,\qquad P(Sm)=V(Sm,m)=W(Sm,m)=m.\qquad}
\]

\subsection{Dependency boundary}

All terms in this section are built from $0,S,I$ by finite substitution.
All inductions have a single equation as their invariant.  Numerical
cases and equality transport are the already reconstructed logical
resources.  Successor injectivity may now be obtained by applying $P$
to an assumed equation $Sx=Sy$.  Equality tests, reflection of
quantifier-free formulas, pairing, and the recovery of general
primitive recursion can then be developed in the established order.
They were not used in order to manufacture predecessor.

This argument is a written internal derivation in the stated theory.
The finite implementation checks do not constitute a proof-assistant
formalization of this entire metatheorem.

% END reference/foundations/scalar_predecessor.tex


% BEGIN reference/foundations/equality_reflection.tex
% Standalone section.  Requires only amsmath and amssymb in the parent document.
% No theorem environments, macros, or bibliography keys are assumed.
% All arithmetic claims in this section are proved here.

\section{From literal-equation induction to full intuitionistic PRA}
\label{sec:eic-completion}\label{app:arithmetic}

This section supplies an unconditional completion theorem. It is adapted
with attribution from the current forward-left draft~\cite{forward}. In the
scalar application used here, every function occurring in the proof is
already a scalar term: predecessor is supplied by
Appendix~\ref{sec:hybrid-bootstrap}, the selector is
$C(n,a,b)=I(n,a;u.b)$, addition is $x+y=I(y,x;u.Su)$, and
truncated subtraction is $d(x,y)=I(y,x;u.Pu)$. All other functions below
are finite compositions of these. Thus this proof does not presuppose
that arbitrary primitive recursion has already been compiled. The QFI
conclusion first applies to the scalar language; the general PR compiler
comes afterward.  Literal-equation
induction, together with the primitive-recursive defining equations,
predecessor coverage, and separation of $0$ and $1$, proves induction for every
quantifier-free formula.  In particular, the proof does not assume a faithful
numerical equality test, stability or decidability of equality, successor
uniqueness, or any commutation law for an iterator.  We construct the equality
test and prove its reflection property for arbitrary elements of the
numerical domain.

\subsection{The exact hypotheses}

Work in ordinary intuitionistic first-order logic with equality and genuine
falsity $\bot$.  The language has one numerical sort and a symbol for each
fixed, finite primitive-recursive description.  Distinct descriptions are
not identified merely because they compute the same standard function.
Let $\mathsf{PRE}$ contain the defining equations for zero, successor,
projections, composition, and primitive recursion.  In particular, choose
descriptions $P$ and $C$ with equations
\begin{align}
 P(0)&=0, & P(Sx)&=x, \label{eq:eic-predecessor}\\
 C(0,a,b)&=a, & C(Sx,a,b)&=b. \label{eq:eic-selector}
\end{align}
Write $1=S0$.  In addition to $\mathsf{PRE}$, assume
\begin{align}
 \mathsf{Cov}:&\qquad
       \forall x\,\bigl(x=0\lor x=S(Px)\bigr),
       \label{eq:eic-coverage}\\
 \mathsf{Sep}:&\qquad 0=1\longrightarrow\bot,
       \label{eq:eic-separation}
\end{align}
and induction for each \emph{literal equation} $E(n,\vec p)$:
\begin{equation}
 \forall\vec p\,
 \left[
   \left(E(0,\vec p)\land
       \forall n\bigl(E(n,\vec p)\longrightarrow E(Sn,\vec p)\bigr)
   \right)
   \longrightarrow\forall n\,E(n,\vec p)
 \right].
 \label{eq:eic-induction}
\end{equation}
Here the schematic expression $E$ is a single equality of numerical
terms; the displayed universal quantifiers range only over numbers.
Every description and every instance of the schema is finite external
syntax.

The axiom formulation \eqref{eq:eic-induction} and the usual contextual
induction rule are interchangeable.  For example, if
\[
 \Gamma\vdash E(0,\vec p),
 \qquad
 \Gamma,E(z,\vec p)\vdash E(Sz,\vec p),
\]
where $z$ is absent from $\Gamma$ and the passive parameters, discharge the
second premise's equation and universally introduce $z$.  The axiom then
gives $\Gamma\vdash\forall n E(n,\vec p)$, which can be instantiated at any
target term.  Conversely, assume the antecedent of the axiom and apply the
contextual rule with fresh variables.  This uses the ordinary intuitionistic
logical rules and permits arbitrary, including quantified, formulas in
$\Gamma$.

\subsection{Constructor separation and two elementary tests}

Define the sign description and the countdown description by
\begin{equation}
 \operatorname{sg}(u)=C(u,0,1),
 \qquad
 K(u)=C(u,1,Pu).
 \label{eq:eic-sign-countdown}
\end{equation}
Their defining equations imply
\begin{align}
 \operatorname{sg}(0)&=0,
 &\operatorname{sg}(Sx)&=1,\label{eq:eic-sign-computation}\\
 K(0)&=1,
 &K(Sx)&=x.\label{eq:eic-countdown-computation}
\end{align}

\paragraph{Constructor laws.}
Successor is injective: applying $P$ to $Sx=Sy$ gives $x=y$ by
\eqref{eq:eic-predecessor}.  Successor never equals zero.  Indeed,
$Sx=0$ gives $1=0$ by applying $\operatorname{sg}$ and using
\eqref{eq:eic-sign-computation}; symmetry and
\eqref{eq:eic-separation} give $\bot$.  Consequently $1\ne0$, where
$A\ne B$ abbreviates $A=B\to\bot$.  None of these arguments uses
induction.

\paragraph{The sign test.}
For every $u$,
\begin{align}
 &\operatorname{sg}(u)=0\lor\operatorname{sg}(u)=1,
       \label{eq:eic-sign-range}\\
 &\operatorname{sg}(u)=0\quad\longleftrightarrow\quad u=0,
       \label{eq:eic-sign-zero}\\
 &\operatorname{sg}(u)=1\quad\longleftrightarrow\quad u\ne0.
       \label{eq:eic-sign-one}
\end{align}
To prove the range assertion, split \eqref{eq:eic-coverage}.  The branch
$u=0$ gives sign $0$, and the branch $u=S(Pu)$ gives sign $1$.
The implication $u=0\to\operatorname{sg}(u)=0$ is computation.
For its converse, the successor branch of coverage would make
$\operatorname{sg}(u)=0$ imply $1=0$, a contradiction; in that branch
explosion gives $u=0$ as well.  Similarly, if
$\operatorname{sg}(u)=1$ and $u=0$, computation gives $0=1$.
Conversely, $u\ne0$ excludes the zero branch of coverage, and the
successor branch computes the sign as $1$.

Coverage and \eqref{eq:eic-selector} also give, for all $u,v$,
\begin{equation}
 C(u,v,v)=v.
 \label{eq:eic-selector-idempotence}
\end{equation}
Both coverage branches compute to $v$, so this identity again needs no
induction.

\paragraph{The countdown test and the notation $W$.}
For every $u$,
\begin{equation}
 K(u)=0\quad\longleftrightarrow\quad u=1.
 \label{eq:eic-countdown-zero}
\end{equation}
If $u=1=S0$, \eqref{eq:eic-countdown-computation} gives $K(u)=0$.
Conversely, assume $K(u)=0$ and split coverage.  If $u=0$, then
$K(u)=1$, contradicting separation, and explosion gives $u=1$.
If $u=S(Pu)$, computation gives $K(u)=Pu$; hence $Pu=0$ and
$u=S0=1$.

If a presentation uses a description named $W$ with defining equations
\[
 W(0)=1,\qquad W(Sx)=x,
\]
then it proves $K(u)=W(u)$ for every $u$ already from coverage.
In the zero branch both sides are $1$; in the branch $u=S(Pu)$ both
sides are $Pu$.  Thus $K$ and $W$ are provably equal on the whole domain,
even if they are distinct descriptions.  This equality does not appeal to
standard-model agreement or to uniqueness of an arbitrary recurrence.

\subsection{Four elementary uses of equation induction}

Choose descriptions for addition and truncated subtraction, with recursion
on the second argument:
\begin{align}
 x+0&=x, &x+Sy&=S(x+y),\label{eq:eic-addition}\\
 d(x,0)&=x, &d(x,Sy)&=P(d(x,y)).\label{eq:eic-subtraction}
\end{align}
Using these equations and \eqref{eq:eic-induction}, we now derive
\begin{align}
 0+y&=y,\label{eq:eic-left-zero}\\
 (Sx)+y&=S(x+y),\label{eq:eic-left-successor}\\
 d(Sx,Sy)&=d(x,y),\label{eq:eic-cancel-successors}\\
 d(Sa,a)&=1.\label{eq:eic-diagonal}
\end{align}
Each induction below has a single equation as its induction formula.

For \eqref{eq:eic-left-zero}, induct on $y$.  The base is $0+0=0$.
The step is $0+Sy=S(0+y)=Sy$.

For \eqref{eq:eic-left-successor}, hold $x$ fixed and induct on $y$.
At zero both sides are $Sx$.  At a successor the induction hypothesis
gives
\[
 (Sx)+Sy=S\bigl((Sx)+y\bigr)
       =S\bigl(S(x+y)\bigr)=S(x+Sy).
\]
Combining \eqref{eq:eic-left-successor} with the defining equation for
addition gives the useful rearrangement
\begin{equation}
 (Sy)+r=y+Sr.
 \label{eq:eic-shift-successor}
\end{equation}
Indeed, both sides are $S(y+r)$.

For \eqref{eq:eic-cancel-successors}, hold $x$ fixed and induct on $y$.
The base computes as
\[
 d(Sx,S0)=P(d(Sx,0))=P(Sx)=x=d(x,0).
\]
Applying $P$ to the induction equation gives the step:
\[
 d(Sx,S(Sy))=P(d(Sx,Sy))=P(d(x,y))=d(x,Sy).
\]

Finally, induct on $a$ for \eqref{eq:eic-diagonal}.  The base is
$d(S0,0)=S0=1$.  At a successor,
\[
 d(S(Sa),Sa)=d(Sa,a)=1,
\]
using \eqref{eq:eic-cancel-successors} and the induction hypothesis.

\subsection{An equation invariant that reflects actual equality}

The remaining arithmetic ingredient is the following equation, valid for
all $x,y$:
\begin{equation}
 C\bigl(\operatorname{sg}(d(x,y)),\ x,\ y+d(x,y)\bigr)=x.
 \label{eq:eic-residual-invariant}
\end{equation}
Its meaning is that a nonzero subtraction residual reconstructs the
original argument by addition.  The selector expresses this assertion as
one literal equation, so equation induction applies directly.

Hold $x$ fixed and induct on $y$.  At $y=0$, subtraction gives residual
$x$, and \eqref{eq:eic-left-zero} gives $0+x=x$.  The desired equation
is consequently $C(\operatorname{sg}(x),x,x)=x$, which is
\eqref{eq:eic-selector-idempotence}.

For the step, assume \eqref{eq:eic-residual-invariant} at $y$ and
abbreviate
\[
 r=d(x,y),\qquad r'=Pr=d(x,Sy).
\]
These are local term abbreviations, not additional function symbols.
Split coverage for $r'$.

If $r'=0$, then $\operatorname{sg}(r')=0$, and the successor instance
of \eqref{eq:eic-residual-invariant} reduces immediately to $x=x$.

In the other branch, $r'=S(Pr')$, so constructor separation gives
$r'\ne0$.  Split coverage for $r$.  Its zero branch would imply
$r'=P0=0$, a contradiction, from which the desired step equation follows.
In its successor branch,
\[
 r=S(Pr)=Sr',\qquad \operatorname{sg}(r)=1.
\]
The induction equation therefore computes to $y+r=x$.  Meanwhile,
\eqref{eq:eic-shift-successor} gives
\[
 (Sy)+r'=y+Sr'=y+r=x.
\]
Since $r'$ is a successor, $\operatorname{sg}(r')=1$, so this is
exactly the branch selected in the successor instance of
\eqref{eq:eic-residual-invariant}.  This completes the equation induction.
All case distinctions came from coverage; no induction on an implication
or a disjunction was used.

\paragraph{Reflection at residual $1$.}
We have, for arbitrary $a,b$,
\begin{equation}
 d(Sa,b)=1\quad\longleftrightarrow\quad a=b.
 \label{eq:eic-residual-reflection}
\end{equation}
For the forward direction, substitute $x=Sa$, $y=b$ into
\eqref{eq:eic-residual-invariant}.  If $d(Sa,b)=1$, the sign is $1$
and the selector returns its second branch, giving
\[
 b+1=Sa.
\]
The defining equation for addition gives $b+1=S(b+0)=Sb$.  Thus
$Sb=Sa$, and applying $P$ gives $b=a$, hence $a=b$.
Conversely, $a=b$ and \eqref{eq:eic-diagonal} give
$d(Sa,b)=d(Sa,a)=1$.

Define the equality characteristic by the fixed description
\begin{equation}
 e(a,b)=\operatorname{sg}\bigl(K(d(Sa,b))\bigr).
 \label{eq:eic-equality-characteristic}
\end{equation}
The sign, countdown, and residual reflection lemmas give
\begin{align}
 e(a,b)=0&\quad\longleftrightarrow\quad a=b,
       \label{eq:eic-equality-reflection}\\
 e(a,b)=0&\ \lor\ e(a,b)=1.
       \label{eq:eic-equality-range}
\end{align}
In detail, $e(a,b)=0$ is equivalent to $K(d(Sa,b))=0$ by
\eqref{eq:eic-sign-zero}, which is equivalent to $d(Sa,b)=1$ by
\eqref{eq:eic-countdown-zero}, and then to $a=b$ by
\eqref{eq:eic-residual-reflection}.  The range assertion is
\eqref{eq:eic-sign-range} applied to $K(d(Sa,b))$.

The reflection property concerns the logic's actual equality throughout
the domain.  It was not assumed as part of a numerical comparison symbol,
and its proof did not restrict $a$ or $b$ to standard numerals.
It also yields decidable equality.  In the $e(a,b)=0$ branch,
\eqref{eq:eic-equality-reflection} gives $a=b$.  In the $e(a,b)=1$
branch, an assumption $a=b$ would give $e(a,b)=0$ and therefore $0=1$,
contradicting separation.  Hence
\begin{equation}
 \forall a\,\forall b\,\bigl(a=b\lor(a=b\to\bot)\bigr)
 \label{eq:eic-equality-decidable}
\end{equation}
is a theorem, not an extra hypothesis.

\subsection{Compilation of all quantifier-free formulas}

For each fixed quantifier-free formula $A(\vec p)$, construct a term
$b_A(\vec p)$ by external recursion on its finite syntax.  Take equality,
$\bot$, $\land$, $\lor$, and $\to$ as the primitive formula constructors;
negation is implication into $\bot$.  Set
\begin{align}
 b_{s=t}&=e(s,t), & b_{\bot}&=1,\label{eq:eic-boolean-atoms}\\
 b_{A\land D}&=C(b_A,b_D,1),\label{eq:eic-boolean-and}\\
 b_{A\lor D}&=C(b_A,0,b_D),\label{eq:eic-boolean-or}\\
 b_{A\to D}&=C(b_A,b_D,0).\label{eq:eic-boolean-implies}
\end{align}
Every term here is a finite composition of fixed primitive-recursive
descriptions; if desired, projections and composition give it a single
description symbol.  No internal quantification over formulas or programs
is involved.

An external induction on the syntax of $A$ proves, within the arithmetic
theory,
\begin{equation}
 \bigl(b_A=0\lor b_A=1\bigr)
 \quad\text{and}\quad
 \bigl(b_A=0\longleftrightarrow A\bigr).
 \label{eq:eic-boolean-correctness}
\end{equation}
Here are the full connective checks.  In every case $b_A=1$ implies
$\neg A$: if $A$ held, the already established implication
$A\to b_A=0$ would contradict separation.

\paragraph{Atomic formulas and falsity.}
For an equation, both claims are
\eqref{eq:eic-equality-reflection}--\eqref{eq:eic-equality-range}.
For falsity the range assertion is $1=0\lor1=1$, with its right disjunct
proved by reflexivity.  Also $1=0\to\bot$ follows from separation and
symmetry, while $\bot\to1=0$ follows by explosion.  Thus $1=0$ is
equivalent to $\bot$, as required.

\paragraph{Conjunction.}
Split $b_A=0\lor b_A=1$.  In the zero branch, $A$ holds and
$b_{A\land D}=b_D$.  Its range follows from the range of $b_D$,
and its being zero is equivalent to $D$, hence to $A\land D$ under
the known $A$.  In the one branch, $\neg A$ holds and
$b_{A\land D}=1$.  Its range is immediate.  If it were zero,
separation would give $\bot$, hence $A\land D$; conversely an
assumption $A\land D$ contradicts $\neg A$, and explosion gives
$b_{A\land D}=0$.

\paragraph{Disjunction.}
In the branch $b_A=0$, the formula $A$ holds and
$b_{A\lor D}=0$.  The range is immediate, and both $b_{A\lor D}=0$
and $A\lor D$ hold, proving their equivalence.  In the branch
$b_A=1$, we have $\neg A$ and $b_{A\lor D}=b_D$.  Its range follows
from the range of $b_D$.  If its value is zero, then $D$ and hence
$A\lor D$.  Conversely, from $A\lor D$, the $A$ alternative
contradicts $\neg A$ and yields $b_D=0$ by explosion, while the
$D$ alternative yields $b_D=0$ directly.  This proves the required
equivalence in both alternatives.

\paragraph{Implication.}
In the branch $b_A=0$, the formula $A$ holds and
$b_{A\to D}=b_D$.  Its range follows from the range of $b_D$.
If its value is zero, then $D$, which gives $A\to D$.  Conversely,
$A\to D$ and the known $A$ give $D$, hence $b_D=0$.
In the branch $b_A=1$, we have $\neg A$ and
$b_{A\to D}=0$.  Its range is immediate.  To prove $A\to D$,
assume $A$; this contradicts $\neg A$, and explosion gives $D$.
Thus both $b_{A\to D}=0$ and $A\to D$ hold in this branch, and
their equivalence follows.  This argument uses the proved two-value
range of the finite characteristic.  It does not use excluded middle
for arbitrary formulas or change the surrounding intuitionistic logic.

This completes the external syntactic induction proving
\eqref{eq:eic-boolean-correctness} for every quantifier-free formula.

\subsection{Recovery of induction and the full theory}

Let $A(n,\vec p)$ be any fixed quantifier-free formula.  Suppose an
arbitrary retained context $\Gamma$ proves its base and its induction
step:
\[
 \Gamma\vdash A(0,\vec p),
 \qquad
 \Gamma,A(z,\vec p)\vdash A(Sz,\vec p),
\]
with the usual freshness condition on $z$.  By
\eqref{eq:eic-boolean-correctness}, the base gives
$\Gamma\vdash b_A(0,\vec p)=0$.  Under the additional assumption
$b_A(z,\vec p)=0$, the same equivalence gives $A(z,\vec p)$,
the given step gives $A(Sz,\vec p)$, and the equivalence gives
$b_A(Sz,\vec p)=0$.  Apply literal-equation induction to
$b_A(n,\vec p)=0$ and then apply its reflection implication.
The result is $\Gamma\vdash A(t,\vec p)$ for every target term $t$.
Thus the theory proves the full contextual quantifier-free induction
rule, including in arbitrary quantified contexts.

\paragraph{Completion theorem.}
In the common language of fixed primitive-recursive descriptions, the
theory
\[
 \mathsf{PRE}+\mathsf{Cov}+\mathsf{Sep}+\mathsf{EI}
\]
has exactly the same sequents as ordinary first-order intuitionistic
primitive recursive arithmetic with all quantifier-free induction
instances.  The forward simulation follows from the recovered
quantifier-free induction, successor injectivity, and successor separation;
all primitive-recursive defining equations and all logical rules are
shared.  For the reverse simulation, literal-equation induction is an
instance of quantifier-free induction, and successor separation gives
$\mathsf{Sep}$.  Quantifier-free induction proves $\mathsf{Cov}$:
at zero choose its left disjunct, and at a successor use
$Sx=S(P(Sx))$, which follows from $P(Sx)=x$, for its right disjunct.

Both interpretations therefore use the full domain, defined by $x=x$,
and preserve every function symbol, equality, logical connective, and
quantifier.  Both composites have the identity isomorphism, with graph
$x=y$.  The theories are consequently identity-domain bi-interpretable
and equiconsistent.  This completion theorem is unconditional once the
literal-equation induction rule has been obtained; it has no dependency
on a classical successor-uniqueness reduction.

% END reference/foundations/equality_reflection.tex


% BEGIN reference/foundations/primitive_recursion.tex
\section{Pairing, primitive recursion, and internal numerical composites}
\label{sec:iterator-coding}
This appendix is adapted with attribution from the progress-closure
arithmetic development~\cite{progress}. Its prerequisites have now been
proved: scalar predecessor, arithmetic reflection, and full quantifier-free
induction. Pairing is a derived proof tool; it is not part of any checker.

\subsection{Pairing from a bounded square-root search}

Once QFI is available, a familiar square-based injection gives a
particularly direct compiler. Write $a\le b$ for $d(a,b)=0$.
The usual linear-order and semiring identities follow by fixed
quantifier-free inductions from the successor, addition, and monus
equations. Multiplication itself is the scalar term
\[
                        a\cdot b=I(b,0;u.u+a).
\]
Only primitive-recursive arithmetic is involved in the following
bounded root search; there is no minimization over an unbounded domain.

For a fixed parameter $n$, put
\[
 U_n(v)=C(d((Sv)^2,n),Sv,v),\qquad
 R_n(i)=I(i,0;v.U_n(v)),\qquad q(n)=R_n(n).
\]
The update increments $v$ exactly when $(v+1)^2\le n$, and
otherwise retains it. QFI in $i$, with $n$ fixed, proves
\[
 R_n(i)\le i,\qquad R_n(i)^2\le n,\qquad
 R_n(i)=i\ \vee\ n<(S(R_n(i)))^2.              \tag{Root invariant}
\]
The base has $R_n(0)=0$. If the update increments, its test
contradicts the strict-upper-bound alternative in the third invariant.
Thus the old state equals $i$, and the new state equals $Si$;
its square bound is the increment test. If the update does not increment, the
strict upper square bound holds and the state remains fixed.
An already established strict upper square bound forces every
subsequent update to remain fixed. These are exactly the cases
needed for the three displayed invariants.

At $i=n$, either the strict upper square bound already holds or
$R_n(n)=n$, in which case $n<(n+1)^2$ supplies it. Thus
\[
                        q(n)^2\le n<(q(n)+1)^2.           \tag{Root}
\]
Monotonicity of squares shows uniqueness of a number satisfying
Root. All formulas just used are quantifier free in the scalar
language, with fixed $n$ during the induction.

Now define
\[
 \pi(a,b)=(a+b)^2+a,\qquad
 L(n)=d(n,q(n)^2),\qquad H(n)=d(q(n),L(n)).
\]
For $k=a+b$, one has $k^2\le k^2+a<(k+1)^2$ since $a\le k$.
Root uniqueness gives $q(\pi(a,b))=a+b$, and therefore
\[
                    L(\pi(a,b))=a,\qquad H(\pi(a,b))=b.
\]
Surjectivity of $\pi$ is unnecessary. The original numerical domain
is never replaced by the range of this coding; pairing is used only
to store the finite state of a fixed primitive recursion. This
standard square pairing is a convenient reconstruction here, not
a claimed new arithmetic basis.

\subsection{Compiling each primitive recursion}

Suppose the fixed target description is
\[
 f(0,\vec p)=g(\vec p),\qquad
 f(Sn,\vec p)=h(n,f(n,\vec p),\vec p).
\]
Assume by external structural induction that the earlier $g,h$
already have scalar translations $g^*,h^*$. Define
\[
 W(u,\vec p)=\pi(SL(u),h^*(L(u),H(u),\vec p)),
\]
\[
 V(n,\vec p)=I(n,\pi(0,g^*(\vec p));u.W(u,\vec p)),
 \qquad f^*(n,\vec p)=H(V(n,\vec p)).
\]
The update has only its state variable bound, with every $p_i$
fixed. Equation induction proves $L(V(n,\vec p))=n$.
The two projection equations and scalar computation then prove
both defining equations for $f^*$ internally. Initial functions
and composition translate directly. This completes the translation
of every fixed finite PR description without adding a universal
evaluator.

\subsection{Why both numerical round trips are internal}

Two fixed scalar iterators with the same seed and provably equal
updates are internally equal at all counters: use literal-equation
induction comparing their values. Their base agrees, and their
successor equations plus congruence prove the step. The same
uniqueness argument works for ordinary primitive recursion.

For an ordinary PR description, induct externally on its finite
construction. The double translation satisfies its original
defining equations, so internal equation induction identifies it
with the original operation. For a native scalar description,
earlier translated updates agree by the external induction
hypothesis; scalar uniqueness identifies the double-translated
iterator. These proofs hold at arbitrary internal counters.

Thus the domain relation $x=y$ is provably a total bijection
preserving every fixed numerical operation in both composites.
The formula constructions and round trips in
Section~\ref{sec:choice-introduction} give the corresponding formula
composites. Both interpretations preserve falsity, taking native
$0=1$ to an ordinary contradiction and ordinary $\bot$ to native
$0=1$. Finite proof translations therefore establish equiconsistency.

% END reference/foundations/primitive_recursion.tex


% BEGIN reference/foundations/inherited_ordinary_bridge.tex
\section{Arithmetic strength and both ordinary interpretation composites}
\label{sec:ordinary-bridge}

We spell out the arithmetic bridge to avoid identifying standard numerical
agreement with bi-interpretability. First elaborate every fixed iterator
into an ordinary function symbol. At stage zero the symbols are $0,S$.
Given a fixed earlier-stage term $h(u,\vec p)$, introduce
$i_h(n,a,\vec p)$ with the two equations
\[
i_h(0,a,\vec p)=a,\qquad
i_h(Sn,a,\vec p)=h(i_h(n,a,\vec p),\vec p).
\]
Only finite descriptions occur; the domain contains no code for a
universal evaluator. Formal parameters are distinct when a description
is declared, while actual arguments can coincide.

Let $T_I$ be ordinary intuitionistic first-order logic with actual
equality and genuine falsity, these staged equations, contextual
induction for literal equations, and
\[
\forall x\,(x=0\vee\exists y\,x=Sy),
\qquad (0=S0)\to\bot.
\]
The second sentence is the falsity bridge. It must be stated: positive
numerical equations by themselves admit a singleton structure. Native
explosion from $0=S0$ and a translation of ordinary falsity to that
equation do not authorize silently dropping the bridge from $T_I$.

The native-to-ordinary translation is
\[
(s=t)^\circ=(s=t),\qquad
(\mathsf D_{x,y}B)^\circ=
\exists x\bigl(B^\circ(x,0)\to\forall y B^\circ(x,y)\bigr).
\]
A clause $\langle A;y.B\rangle$ becomes $A^\circ\to\forall yB^\circ$.
For the auxiliary closure used by the compact checkers,
$(\mathsf Q_yB)^\circ=B^\circ(0)\to\forall yB^\circ$.
The interpretation domain is all elements, defined by $x=x$;
equality is actual equality.

Case DI is sound because numerical coverage combines its retained base
with its successor proof. Progress DI is sound by literal-equation
induction under a temporarily assumed base. Use, DE and transport are
ordinary logical inferences under the translations. Conversely the
derived logical operations, equation induction, cases and explosion in
Section~\ref{sec:choice-introduction} recover the rules and axioms of
$T_I$. An auxiliary closure is interderivable with
$\mathsf D_{v,y}B(y)$ for unused $v$: in one direction use case DI with
the closure and its base; in the other eliminate the paired formula
and return its supplied closure. Thus formulas and context clauses have
the stated meanings in both presentations.

\subsection{The arithmetic dependency order}

The scalar update has no recursion-counter argument. We do not conceal
one in an encoding. The inherited internal proofs, reproduced in the
appendices, establish the required arithmetic in this order:
\begin{enumerate}
\item Using equation induction and numerical cases, construct a fixed
scalar predecessor term and prove $P0=0$ and $P(Sn)=n$.
Existential coverage then gives $x=0\vee x=S(Px)$: in the successor
case $x=Sy$, apply $P$ to obtain $Px=y$ and substitute.
\item Construct truncated subtraction and a numerical characteristic
of actual equality. Prove its reflection on the entire original domain.
\item Compile each fixed quantifier-free formula into its characteristic
equation and recover every contextual quantifier-free induction instance.
\item With that induction available, construct pairing and projections
by a bounded square-root search and compile every fixed general
primitive-recursive description into a scalar iterator term.
\end{enumerate}
No pairing or faithful equality detector is assumed to obtain
predecessor. No appeal to standard-model correctness replaces an
internal equality. Every induction fixes its passive parameters;
diagonal substitution occurs only into a previously proved identity.

For a target recursion
\[
f(0,\vec p)=g(\vec p),\qquad
f(Sn,\vec p)=h(n,f(n,\vec p),\vec p),
\]
let $\pi,L,H$ be the proved pairing and projections. The compiler uses
\[
V(n,\vec p)=I\bigl(n,\pi(0,g^*(\vec p));
u.\pi(SL(u),h^*(L(u),H(u),\vec p))\bigr),
\qquad f^*(n,\vec p)=H(V(n,\vec p)).
\]
Equation induction proves $L(V(n,\vec p))=n$; the two projection
identities then prove the target defining equations. The construction
is external recursion on each finite PR description, not an internal
quantification over arbitrary functions.

\begin{theorem}[Whole-domain bi-interpretability and equiconsistency]
The ordinary realization $T_I$ and full first-order iPRA are
bi-interpretable. In both theories $x=y$ defines the isomorphism from
the original domain to its double interpretation. The native Choice-DI
calculus is therefore a faithful proof presentation of a theory
bi-interpretable and equiconsistent with iPRA.
\end{theorem}
\begin{proof}
Interpret an iterator in iPRA by primitive recursion that ignores its
counter argument in the update. iPRA proves its equations, literal
equation induction, coverage and separation. In the reverse direction
use the scalar compiler. The predecessor and equality-reflection proofs
supply the successor axioms and every quantifier-free induction instance.
All ordinary logical rules are preserved.

For two iterators with the same seed and provably equal updates,
equation induction on the equality of their values proves uniqueness at
every counter, with all parameters fixed. The corresponding argument
applies to ordinary primitive recursion. Induct externally on a finite
description: its double translation satisfies its original defining
equations, so internal uniqueness identifies the two functions. This
also proves coherence when substitution changes the staged syntactic
description of an update.

The domain and equality are unchanged; hence the identity graph is
internally total, bijective and preserves every interpreted operation.
The displayed ordinary meanings of the logical macros and the paired
binder give both formula composites by external induction on finite
formula syntax. Clauses are similarly equivalent to their conditional
universal meanings. The explicit falsity bridge preserves contradiction
in both directions. Effective translations of finite proofs give
equiconsistency, without a claim that either theory proves its own
consistency.
\end{proof}

% END reference/foundations/inherited_ordinary_bridge.tex

\clearpage
\part{The source restriction used by the compact checker}

% BEGIN reference/zero_source_admissibility.tex
% Standalone mathematical module; requires amsmath, amssymb, amsthm, url.
% The surrounding document should supply theorem, lemma, and corollary environments.
% Written derivation, not a proof-assistant verification.
\section{Removing literal zero from the source of transport}
\label{zs:section}

The following is a restriction of a mathematical inference rule, rather
than a change in the spelling of zero. Throughout this section,
``nonzero source'' means \emph{not syntactically the single term $0$}.
Every variable and every term headed by $S$ or $I$ is a permitted source,
including terms whose numerical value is zero.

Let $\mathcal C$ be the native scalar calculus with terms
\[
 t::=x\mid0\mid St\mid I(n,a;u.h),
\]
the two computation schemas
\[
 I(0,a;u.h)=a,\qquad
 I(Sn,a;u.h)=h[I(n,a;u.h)/u],
\]
formulas generated by equality and the $Q,D$ binders, and the Hyp, QE,
DE, case-DI and equation-progress-DI rules stated in the main inventory.
The same argument applies to the original presentation with $D$ object
formulas and $Q$ resources in contexts. Let $\mathcal C^{-}$ replace
unrestricted transport by
\[
 \frac{\Gamma\vdash s=t\qquad\Gamma\vdash F(s)}
      {\Gamma\vdash F(t)}
 \quad\mathsf{Tr}^{-},\qquad s\not\equiv0.
 \tag{Tr$^-$}\label{zs:restricted-rule}
\]
Here $F$ is an arbitrary native numerical template, and $s\not\equiv0$
is a condition on finite term syntax. All other rules are unchanged.

\begin{theorem}[Literal-zero source elimination]
\label{zs:main}
For every native context $\Gamma$ and formula $A$,
\[
 \Gamma\vdash_{\mathcal C^{-}}A
 \quad\Longleftrightarrow\quad
 \Gamma\vdash_{\mathcal C}A.
\]
In particular, the ordinary whole-domain realization of
$\mathcal C^{-}$ is bi-interpretable and equiconsistent with full
first-order iPRA through the same interpretations as the native scalar
core. No change of numerical domain or interpreted equality is needed.
\end{theorem}

We prove the theorem without assuming arbitrary numerical substitution
admissibility for $\mathcal C^{-}$: such substitution can turn a variable
source into the forbidden literal zero. Every occurrence of a helper
equation at an arbitrary term is either obtained directly from its
displayed derived rule at that term, or from a derivation whose transport
sources retain an $S$ or $I$ head under the indicated substitution.

\subsection{Structural operations, cases, and equation induction}

Weakening, ordinary-assumption cut, and capture-free renaming are
admissible by induction on derivations. These operations do not turn a
nonzero source into literal zero. As usual, fresh parameters are renamed
before a cut. We also use the following restricted substitution fact:
a capture-free substitution is admissible for a particular finite
derivation if it leaves every primitive transport source in that
derivation syntactically different from zero. This fact follows by the
same induction; it is not unrestricted substitution admissibility.

\begin{lemma}[Cases and equation induction need no transport]
\label{zs:cases-induction}
In $\mathcal C^{-}$ the following rules are derivable:
\[
 \frac{\Gamma\vdash A(0)\qquad\Gamma\vdash A(Sz)}
      {\Gamma\vdash A(t)}\quad\mathsf{NC},
 \qquad
 \frac{\Gamma\vdash E(0)\qquad\Gamma,E(z)\vdash E(Sz)}
      {\Gamma\vdash E(t)}\quad\mathsf{EI}.
\]
In NC, $A$ is any native formula; in EI, $E$ is one literal equation.
The parameter $z$ is fresh for $\Gamma$, $t$, and the fixed parameters
of the matrix. Other parameters stay fixed throughout EI.
\end{lemma}
\begin{proof}
Choose an unused witness coordinate $x$. In NC, weaken the successor
proof by $A(0)$ and apply case-DI, with witness zero, to obtain
$D_{x,y}A(y)$. DE opens $Q_yA(y)$, because $x$ is unused. QE with the
given base proof and target $t$ yields $A(t)$, and DE closes.
For EI use progress-DI instead, with the given equation step, and use
the same DE/QE construction. No transport rule occurs.
\end{proof}

These rules will be applied to unconditional helper equations before
they are weakened into the context containing $0=t$. We never use NC
with a case parameter occurring in its retained context.

\subsection{The equality operations available immediately}

\begin{lemma}[Reflexivity and safe equality operations]
\label{zs:safe-equality}
Reflexivity $r=r$ is derivable for every term $r$. From $p=q$ with
$p\not\equiv0$, one can derive $q=p$ and $H(p)=H(q)$ for any numerical
template $H$. Every direct computation can be used in either direction
as an equation; using the reversed equation as a transport major is
permitted only when its new source is syntactically nonzero.
\end{lemma}
\begin{proof}
Put $b=I(0,r;u.0)$. The computation $b=r$, used as both major and
minor of transport with template $v=r$, gives $r=r$. Its source is
the nonempty iterator term $b$.
For symmetry of $p=q$, transport the reflexivity $p=p$ through the
template $v=p$. For numerical congruence, transport the reflexivity
$H(p)=H(p)$ through $H(p)=H(v)$. Both use the permitted source $p$.
Every computation equation has an $I$-headed left side, so its
symmetry follows by the same construction.
\end{proof}

An occurrence-selective template may keep some occurrences of a term
fixed while abstracting other occurrences. This is ordinary
capture-free template substitution, not an extra inference rule.

\begin{lemma}[Joining two nonzero terms known to be zero]
\label{zs:zero-join}
Suppose $p,q\not\equiv0$ and $\Gamma$ proves $p=0$ and $q=0$.
Then $\Gamma\vdash p=q$ using only $\mathsf{Tr}^{-}$ and computation.
\end{lemma}
\begin{proof}
Let
\[
 z_0=I(0,0;u.0),\qquad b_p=I(0,p;u.0).
\]
Computation gives $b_p=p$. Start with $b_p=b_p$ and use the major
$p=0$ in the template $I(0,v;u.0)=b_p$. This gives $z_0=b_p$.
Symmetry at the nonzero source $z_0$ gives $b_p=z_0$.
Now transport through $v=z_0$ using the computation $b_p=p$, whose
source is $b_p$, to obtain $p=z_0$.
The identical construction gives $q=z_0$. Symmetry at source $q$
gives $z_0=q$, and transport at source $z_0$ in $p=z_0$ gives $p=q$.
The sources used are $p,q,z_0,b_p,b_q$; all are syntactically nonzero.
No reverse transport from zero has been used.
\end{proof}

The zero-join lemma is deliberately placed before the next construction.
The positive case of its singleton test needs precisely this lemma;
ordinary transitivity through zero is not yet available.

\subsection{A singleton-one test without predecessor}

Define the finite scalar terms
\[
\begin{aligned}
 C(v,a,b)&:=I(v,a;u.b),& Z(v)&:=C(v,1,0),\\
 K(v)&:=C(v,0,1),& H(v)&:=C(v,1,2),\\
 N(n)&:=I(n,0;u.Z(u)),& R(n)&:=I(n,0;u.H(u)),\\
 L(n)&:=N(R(n)).
\end{aligned}
\tag{Tests}\label{zs:tests}
\]
The update variable in $C$ is unused. Here $1=S0$ and $2=S1$.
The standard-model description of $R$ is $0,1,2,2,\ldots$; the proof
below establishes the needed equations internally and does not assume
this description.

\begin{lemma}[Uniform test equations]
\label{zs:singleton}
The following equations have derivations in $\mathcal C^{-}$,
uniformly in the displayed argument:
\[
 H(v)=S K(v),\qquad R(1)=1,\qquad R(SSv)=2,
\]
\[
 N(1)=1,\quad N(2)=0,\quad
 L(1)=1,\quad L(SSv)=0,\quad
 \boxed{L(Sv)=Z(v)}.
 \tag{Singleton}\label{zs:singleton-equations}
\]
All transport sources in these helper derivations have an $I$ or $S$
head, except for the explicit nonzero endpoints in the zero-join
lemma; those endpoints also have an $I$ head in this application.
\end{lemma}
\begin{proof}
For $H(v)=S K(v)$ use NC. At zero both sides reduce to $1$; at a
successor both reduce to $2$. To join two equations $p=j$, $q=j$
with $j=1$ or $j=2$, reverse $q=j$ at its nonzero source $q$, then
transport at the nonzero numeral $j$. Each reduction of these
particular terms is a forward replacement of an iterator computation.
This proves the helper at an arbitrary target $v$ by NC, without
substituting zero into a variable-source transport.

The equations $N(1)=1$ and $N(2)=0$ are finite computations:
$N(0)=0$, $N(1)=Z(N(0))=1$, and $N(2)=Z(N(1))=0$.
Every replacement source in these calculations is an iterator term.
Likewise $R(1)=H(R(0))=H(0)=1$.

Computation and the first helper give, for any $v$,
\[
 R(Sv)=H(R(v))=S K(R(v)).
\]
Use this equation at the nonzero source $R(Sv)$ in
$R(SSv)=H(R(Sv))$. The resulting right side is
$H(S K(R(v)))=2$ by the successor computation of $C$.
Thus $R(SSv)=2$. These substitutions are all at $I$-headed sources.
Substitute $R(1)=1$ and $R(SSv)=2$ inside $N$ and use the two
computed values of $N$ to obtain $L(1)=1$ and $L(SSv)=0$.

Finally apply NC to the matrix $L(Sv)=Z(v)$.
At zero, $L(1)=1$ and $Z(0)=1$ are joined through the nonzero numeral
$1$. At a successor, $L(SSv)=0$ and $Z(Sv)=0$ are joined by
Lemma~\ref{zs:zero-join}; both left sides are nonempty iterator terms.
This proves the boxed equation at every target. Its use is therefore
an instance of a displayed derived rule, not an appeal to unrestricted
substitution admissibility.
\end{proof}

\begin{lemma}[An unconditional zero-filter identity]
\label{zs:filter}
For every term $t$,
\[
                 t=C(Z(t),t,0)
                 \tag{Filter}\label{zs:filter-equation}
\]
is derivable in the empty context of $\mathcal C^{-}$.
\end{lemma}
\begin{proof}
Apply NC to $v=C(Z(v),v,0)$, with the target $t$.
At zero, forward computations give
$C(Z(0),0,0)=C(1,0,0)=0$; reverse this equation at its nonzero
left side. At $Sv$, forward computations give
$C(Z(Sv),Sv,0)=C(0,Sv,0)=Sv$; again reverse at the nonzero
iterator left side. Thus the base and successor proofs have no
retained assumption depending on the case variable. All numerical
arguments of the final NC target are permitted. In particular, this
identity is proved before any assumption $0=t$ is used.
\end{proof}

\subsection{Turning a zero-source equality around}

\begin{lemma}[Zero-source symmetry]
\label{zs:zero-symmetry}
For every term $t$,
\[
                     0=t\ \vdash_{\mathcal C^{-}}\ t=0.
                     \tag{ZS}\label{zs:zero-symmetry-equation}
\]
\end{lemma}
\begin{proof}
If $t\equiv0$, use reflexivity. Otherwise $t$ is a permitted source.
Work in the fixed context $\Gamma=\{0=t\}$ and put
\[
                  A(n):=I(n,t;u.Su).
\]
We first prove $n=A(n)$ by EI, keeping $t$ fixed.

For the base, computation gives $A(0)=t$. Its symmetry gives
$t=A(0)$. Since $t\not\equiv0$, transport at source $t$ in the
retained assumption $0=t$ gives $0=A(0)$.

For the step, retain $z=A(z)$ with $z$ fresh. This source is a
variable, hence is permitted. Numerical successor congruence from
Lemma~\ref{zs:safe-equality} gives $Sz=S A(z)$.
Computation gives $A(Sz)=S A(z)$, and its symmetry gives
$S A(z)=A(Sz)$. Its source $S A(z)$ is also permitted, so transport
in $Sz=S A(z)$ gives $Sz=A(Sz)$. This is exactly the EI step.
The induction hypothesis is used at the same fixed parameter $t$.

At target $1$, EI gives $1=A(1)$. Forward computation gives
$A(1)=St$: the successor equation produces $S A(0)$, and the
nonzero-source computation $A(0)=t$ is used inside $S$.
Transport at source $A(1)$ consequently gives
\[
                              1=St.
\]
Its source $1$ is permitted. Numerical congruence for $L$ gives
$L(1)=L(St)$. Use $L(1)=1$ at source $L(1)$, and
$L(St)=Z(t)$ at source $L(St)$, obtaining $1=Z(t)$.
Symmetry at source $1$ gives $Z(t)=1$.

Weaken the unconditional Filter identity into $\Gamma$.
Transport at the iterator source $Z(t)$ changes its right side to
$C(1,t,0)$. The latter computes to zero, and its source is again
an iterator term. The resulting equation is $t=0$, as required.
\end{proof}

This proof uses cases only for unconditional helper equations. The
sole induction with the retained assumption $0=t$ has the literal
equation $n=A(n)$ as invariant. It does not perform a forbidden
case distinction while silently changing the retained context.

\subsection{Recovering equality algebra and formula transport}

\begin{lemma}[Full numerical equality algebra]
\label{zs:full-equality}
The calculus $\mathcal C^{-}$ admits unrestricted symmetry,
numerical congruence, and transitivity. In particular it admits
root replacement
\[
 \frac{\Gamma\vdash s=t\qquad\Gamma\vdash s=r}
      {\Gamma\vdash t=r}.
\]
\end{lemma}
\begin{proof}
Symmetry uses Lemma~\ref{zs:safe-equality} when the source is not
zero and Lemma~\ref{zs:zero-symmetry} otherwise.

For numerical congruence from $s=t$, a nonzero $s$ is handled by
Lemma~\ref{zs:safe-equality}. If $s\equiv t\equiv0$, use
reflexivity. If $s\equiv0$ and $t\not\equiv0$, first derive $t=0$,
apply permitted numerical congruence at source $t$ to obtain
$H(t)=H(0)$, and use the now established full symmetry.

For transitivity from $r=s$ and $s=t$, a nonzero middle term $s$
is handled by one permitted transport. If $s\equiv0$ and either
endpoint is literally zero, a given premise is already the
conclusion. Otherwise $r,t\not\equiv0$; reverse $0=t$ to get
$t=0$ and apply the zero-join lemma to $r=0$ and $t=0$.
Root replacement follows by symmetry and transitivity.
These constructions are acyclic: zero-source symmetry was proved
before any unrestricted numerical congruence or transitivity.
\end{proof}

\begin{lemma}[Full transport]
\label{zs:full-transport}
For every native numerical template $F(v)$, the rule
\[
 \frac{\Gamma\vdash s=t\qquad\Gamma\vdash F(s)}
      {\Gamma\vdash F(t)}
\]
is admissible in $\mathcal C^{-}$, without a source restriction.
\end{lemma}
\begin{proof}
Use external induction on the finite number of formula binders in
$F$, uniformly in retained contexts and fixed numerical parameters.
For an equation $F(v)\equiv A(v)=B(v)$, use unrestricted numerical
congruence to obtain $A(s)=A(t)$ and $B(s)=B(t)$. Join these with
$A(s)=B(s)$ using the established equality algebra.

For $F(v)=D_{x,y}B(v,x,y)$, eliminate $F(s)$ at a fresh $a$,
retaining $Q_yB(s,a,y)$. Introduce the desired $F(t)$ by case-DI
with witness $a$. Its local assumption is $B(t,a,0)$.
Full symmetry gives $t=s$, and the external induction hypothesis
for the smaller matrix transports the assumption to $B(s,a,0)$.
QE gives $B(s,a,Sz)$, and the induction hypothesis in the forward
direction gives $B(t,a,Sz)$. Close case-DI and DE. All new
parameters can be chosen fresh for the source equality and the
retained context.

If $Q$ is first-class formula syntax, first observe that its case
introduction is derived without transport: introduce
$D_{u,y}B(y)$ with unused $u$ by case-DI and eliminate it with the
one-node continuation $Q_yB(y)$. For $F(v)=Q_yB(v,y)$, use this
derived introduction for the target. Convert its base assumption
backwards by the induction hypothesis, apply the old Q resource,
and convert the successor conclusion forwards. This is the same
smaller-matrix argument and uses no new numerical induction.
\end{proof}

\begin{proof}[Proof of Theorem~\ref{zs:main}]
Every restricted derivation is an unrestricted one. Conversely,
replace each unrestricted transport node by the derivation supplied
by Lemma~\ref{zs:full-transport}, recursively translating its
premises first. Other inference nodes are unchanged. The result has
the same context and conclusion. Thus the full native consequence
relation, not just closed theorems, is preserved. The existing
native/ordinary and ordinary/iPRA interpretations and both identity
composites therefore transfer unchanged.
\end{proof}

\subsection{Compatibility with restricted computation certificates}
\label{zs:certificate-computation}

Some compact checkers require the initializer of a computation packet
to have nonempty syntax, and encode a successor counter by a nonempty
predecessor. These are certificate restrictions, not restrictions
on the two mathematical computation schemas. They do not introduce
a circular dependence on zero-source transport in the preceding proof.

Use the inherited certified-zero term
\[
                         c_0=I(2,1;u.0).
\]
Its equation $c_0=0$ is a directly admitted successor computation:
the predecessor and initializer are both $1$. For an arbitrary
computation instance with literal-zero initializer or literal-zero
successor predecessor, form its whole equation as a template $J(v)$,
replacing just these offending occurrences by $v$. The instance
$J(c_0)$ is a directly admitted computation. Transport through this
equation template using $c_0=0$ yields $J(0)$, the desired original
instance. Both computation premises are admitted directly, and the
only transport source is the nonempty term $c_0$.

This is exactly the inherited three-node normalization in
\path{Research/certified_zero/tests/producer.py}, function
\texttt{expand\_computation}: one transport and two computation
premises. Its source is $c_0$, even when both defects are corrected
simultaneously. Passive zero occurrences in an update are left alone.
Thus all computation instances used in the mathematical proof can
be expanded before compiling to the source-restricted checker.

\subsection{Scope of the result}

The new statement is the admissibility of all omitted literal-zero
source transports. The native term/formula language, arithmetic
schemas, and target theory remain the inherited ones. This theorem
justifies replacing a two-way equality/Q major-pattern test by a
source-discriminated test in a suitable certificate protocol. It
does not assert a minimum proof-checker size, a shortest proof
expansion, or a new characterization of the primitive recursive
functions. The displayed proof is a written universal derivation;
executable certificates, when supplied, validate particular
instances rather than substitute for this argument.


% END reference/zero_source_admissibility.tex

\clearpage
\part{The inherited progressive orbit calculus and its full round trips}

% BEGIN reference/progressive/new_theory.tex
\providecommand{\llbracket}{[\![}
\providecommand{\rrbracket}{]\!]}
\section{A seven-rule calculus with complete progressive elimination}
\label{sec:new-calculus}

The new presentation keeps the scalar iterator and the three native formula
constructors.  It changes the meaning of an atomic closure: an atomic closure
records a progressive equation, while a compound closure retains the old
conditional-universal meaning.  A single elimination rule recovers both
one-step progress and induction from an atomic closure.  Consequently every
formula of the new language has a faithful ordinary reading; no restricted
public grammar or auxiliary-only formula class is needed.

This is a structural simplification, not a claim of an absolute minimum.  The
current Lean kernel has eight rule constructors.  The proposed kernel has
seven, and its introduction rule commutes directly with numerical
substitution.  The earlier seven-rule \texttt{WPrf} presentation also has seven
constructors, but its introduction distinguishes the literal witness $0$ from
other witnesses.  Here the distinction depends only on whether a matrix is a
literal equation, which numerical substitution cannot change.  The concrete
checker measurements, including comparisons in which this calculus is larger,
are reported separately.\cite{native,witness}

\subsection{Complete primitive specification}
\label{subsec:new-specification}

There is one numerical sort.  Terms are
\[
 t ::= x\mid 0\mid St\mid I(n,a;u.h).
\]
Only $u$ is bound in $h$; other occurrences are free parameters.  Substitution
is capture-free.  Put
\begin{equation}
 1:=S0,\qquad
 \mathsf A(a,n):=I(a,n;u.Su).
 \label{eq:new-add}
\end{equation}
Thus the first argument of $\mathsf A$ is the iteration counter and the
second is its initializer.  This argument order is essential below.

Formulas are
\begin{equation}
 F ::= (s=t)\mid Q_yF\mid D_{x,y}F.
 \label{eq:new-formulas}
\end{equation}
The binder $Q$ binds $y$, and $D$ binds the distinct variables $x,y$.
A \emph{literal equation} has outermost constructor $=$.  A \emph{compound}
formula has outermost constructor $Q$ or $D$.  Define the syntactic selector
\begin{equation}
 \eta_B(a):=
 \begin{cases}
 a,&B\text{ is a literal equation},\\
 0,&B\text{ is compound}.
 \end{cases}
 \label{eq:new-eta}
\end{equation}
Write $\Gamma\vdash_{\mathcal R}F$ for derivability.  Contexts are finite
reusable lists of formulas.  The complete primitive inventory is the following
seven rule schemas:
\[
 \frac{F\in\Gamma}{\Gamma\vdash_{\mathcal R}F}
 \quad\mathrm{Hyp},
\]
\[
 \frac{}{\Gamma\vdash_{\mathcal R}I(0,a;u.h)=a}
 \quad\mathrm{I0},\qquad
 \frac{}{\Gamma\vdash_{\mathcal R}
 I(Sn,a;u.h)=h[I(n,a;u.h)/u]}
 \quad\mathrm{IS},
\]
\begin{equation}
 \frac{\Gamma\vdash_{\mathcal R}s=t\qquad
       \Gamma\vdash_{\mathcal R}p(s)=q(s)}
      {\Gamma\vdash_{\mathcal R}p(t)=q(t)}
 \quad\mathrm{Tr}_{=},
 \label{eq:new-tr}
\end{equation}
\begin{equation}
 \frac{\Gamma\vdash_{\mathcal R}Q_yB(y)\qquad
       \Gamma\vdash_{\mathcal R}B(\eta_B(a))}
      {\Gamma\vdash_{\mathcal R}
         B(\mathsf A(\eta_B(a),n))}
 \quad\mathrm{QE}_{\mathsf A},
 \label{eq:new-qe}
\end{equation}
\begin{equation}
 \frac{\Gamma\vdash_{\mathcal R}D_{x,y}B(x,y)\qquad
       \Gamma,Q_yB(b,y)\vdash_{\mathcal R}F}
      {\Gamma\vdash_{\mathcal R}F}
 \quad\mathrm{DE},
 \label{eq:new-de}
\end{equation}
\begin{equation}
 \frac{\Gamma,B(t,\eta_B(z))\vdash_{\mathcal R}B(t,Sz)}
      {\Gamma\vdash_{\mathcal R}D_{x,y}B(x,y)}
 \quad\mathrm{DI}_{\eta}.
 \label{eq:new-di}
\end{equation}
In \eqref{eq:new-tr}, $p(v),q(v)$ are numerical term templates.  In
\eqref{eq:new-qe}, $a,n$ are arbitrary terms in the conclusion's context.
In DE, $b$ is fresh for $\Gamma,F$ and all fixed parameters of $B$.
In DI, $z$ is fresh for $\Gamma,t$ and all fixed parameters of $B$.
The displayed local assumptions are discharged.  The designated
contradiction is $0=1$.

There are no other primitive axioms, inference rules, induction schemas,
logical connectives, successor axioms, or numerical case principles.
In particular, equality transport in a compound formula is not primitive.

\begin{lemma}[Structural rules and substitution]
\label{lem:new-structural}
Weakening, exchange, contraction, renaming, assumption cut, and capture-free
numerical substitution are admissible in $\mathcal R$.
\end{lemma}
\begin{proof}
Use induction on a finite derivation, renaming eigenparameters before a
substitution or a cut.  A cut replaces each relevant hypothesis leaf by the
supplied derivation, weakened into its local context.  The only new
substitution check is
\[
 (\eta_B(a))\sigma=\eta_{B\sigma}(a\sigma).
\]
It holds because numerical substitution preserves the outermost formula
constructor.  The iterator abbreviation \eqref{eq:new-add} also commutes with
substitution.  Hence QE and DI translate to instances of the same rules;
there is no witness-literal exception to repair.
\end{proof}

\subsection{Equality and the induction bootstrap}
\label{subsec:new-bootstrap}

A useful derived rule is
\begin{equation}
 \frac{\Gamma,B(\eta_B(z))\vdash_{\mathcal R}B(Sz)}
      {\Gamma\vdash_{\mathcal R}Q_yB(y)}
 \quad\mathrm{QI}_{\eta}.
 \label{eq:new-qi}
\end{equation}
To obtain it, apply DI to $D_{x,y}B(y)$ with $x$ unused and witness $0$.
Apply DE to this result; its local hypothesis is precisely $Q_yB(y)$, which
can be exported by Hyp because it is independent of the fresh witness.

\begin{lemma}[Numerical equality and contextual equation induction]
\label{lem:new-ei}
The usual reflexivity, symmetry, transitivity, and numerical congruence rules
are derivable.  For every literal equation $E(y)$,
\begin{equation}
 \frac{\Gamma\vdash_{\mathcal R}E(0)\qquad
       \Gamma,E(z)\vdash_{\mathcal R}E(Sz)}
      {\Gamma\vdash_{\mathcal R}E(t)}
 \quad(z\text{ fresh})
 \label{eq:new-ei}
\end{equation}
is admissible with arbitrary retained side assumptions $\Gamma$.
\end{lemma}
\begin{proof}
Let $J(a)=I(0,a;u.0)$.  I0 gives $J(a)=a$; transport this equation along
itself with template $v=a$ to obtain $a=a$.  The usual equational templates
then give symmetry, transitivity, and numerical congruence.

For induction, the step premise and \eqref{eq:new-qi} give $Q_yE(y)$.
Use \eqref{eq:new-qe} with source $a=0$ and initializer $n=t$.  Its conclusion
is $E(I(0,t;u.Su))$.  I0 and purely equational transport give $E(t)$.
No identity $I(t,0;u.Su)=t$ has been assumed.  Using the counter $t$ and
initializer $0$ in the elimination rule instead would not justify this
bootstrap.
\end{proof}

\begin{lemma}[Elementary iterator identities]
\label{lem:new-addition}
The new calculus proves, with arbitrary parameters,
\begin{equation}
 \begin{aligned}
 \mathsf A(0,n)&=n,&
 \mathsf A(Sa,n)&=S\mathsf A(a,n),\\
 \mathsf A(a,0)&=a,&
 \mathsf A(a,Sn)&=S\mathsf A(a,n),\\
 \mathsf A(a,n)&=\mathsf A(n,a),&
 \mathsf A(a,1)&=Sa.
 \end{aligned}
 \label{eq:new-add-identities}
\end{equation}
\end{lemma}
\begin{proof}
The first row consists of computation instances.  Induct on $a$ for
$\mathsf A(a,0)=a$.  Induct on $a$ again for the initializer-successor
identity: its base is $Sn=Sn$, and its step follows by applying $S$ to the
induction hypothesis.  For commutativity, induct on $a$.  The base is
$n=\mathsf A(n,0)$; the step is
\[
 \mathsf A(Sa,n)=S\mathsf A(a,n)
               =S\mathsf A(n,a)=\mathsf A(n,Sa).
\]
Finally $\mathsf A(a,1)=S\mathsf A(a,0)=Sa$.
All inductions are instances of \eqref{eq:new-ei} on literal equations.
\end{proof}

An atomic closure therefore has the derived one-step elimination
\begin{equation}
 Q_yE(y),\ E(a)\ \vdash_{\mathcal R}\ E(Sa).
 \label{eq:new-step}
\end{equation}
Take initializer $1$ in QE and use \eqref{eq:new-add-identities}.
For a compound matrix, QE at source $0$ and initializer $t$ has the ordinary
conditional-closure effect $Q_yB,B(0)\vdash B(t)$, once transport through
$B$ has been derived.

\begin{lemma}[Full transport and homogeneous binder congruence]
\label{lem:new-transport}
For every formula template $F(v)$,
\begin{equation}
 s=t,\ F(s)\ \vdash_{\mathcal R}\ F(t).
 \label{eq:new-full-tr}
\end{equation}
Moreover, if $B$ and $C$ are uniformly interderivable and have the same root
class (both equations or both compound), their $Q$ closures are
interderivable, and their $D$ closures are interderivable.
\end{lemma}
\begin{proof}
For transport, induct on $F$.  Equations use the primitive rule.  For
$Q_yE(v,y)$ with atomic $E$, introduce the target closure by QI.
Reverse-transport its local $E(t,z)$ to $E(s,z)$, use
\eqref{eq:new-step} for the source closure, and transport the resulting
$E(s,Sz)$ to $E(t,Sz)$.

For $Q_yB(v,y)$ with compound $B$, introduce the target closure from
$B(t,0)$.  The induction hypothesis converts this to $B(s,0)$.
QE for the source closure at initializer $Sz$ gives
$B(s,I(0,Sz;u.Su))$.  The induction hypothesis applied inside the strictly
smaller matrix $B$ removes the I0 redex and changes the parameter $s$ to $t$.
This proves the QI target $B(t,Sz)$.

For $D_{x,y}B(v,x,y)$, eliminate the source binder at a fresh witness $b$.
Introduce the target at that witness.  Convert its local assumption
$B(t,b,\eta_B(z))$ to $B(s,b,\eta_B(z))$ by the induction hypothesis.
At an atomic matrix use \eqref{eq:new-step}; at a compound matrix use QE
with source $0$, initializer $Sz$, and I0.  Transport the successor result
back to the parameter $t$.  Close DI and DE.  Every invoked formula
transport concerns the strictly smaller matrix, so this is not circular.

For homogeneous congruence, use the same constructions, replacing numerical
transport by the two given conversions between $B$ and $C$.  Their root
classes agree, so their local DI assumptions use the same value of $\eta$.
For $D$, eliminate the source at a fresh witness and introduce the target
at that witness.  The argument works uniformly under arbitrary additional
retained assumptions.
\end{proof}

The root-class qualification is substantive.  The new $Q$ is not claimed to
preserve arbitrary equivalence when one matrix is an equation and the other
is compound.  A progressive assertion and a conditional-universal assertion
are different assertions even when their matrices are pointwise equivalent.
The translations below respect this boundary explicitly.

\subsection{Ordinary soundness}
\label{subsec:new-soundness}

Use the pinned repository's scalar theory $T_{\rm eq}$: ordinary
intuitionistic first-order logic with actual equality, the two iterator
computations, numerical coverage
\[
 \forall x\,(x=0\lor\exists y\,x=Sy),
\]
separation $(0=1)\to\bot$, and contextual induction for literal equations.
The repository proves its whole-domain bi-interpretability with iPRA.\cite{bridge}

Define the new ordinary reading recursively by
\begin{equation}
 \begin{aligned}
 \llbracket s=t\rrbracket&=(s=t),\\
 \llbracket Q_yE(y)\rrbracket
   &=\forall y\,(E(y)\to E(Sy))
       &&(E\text{ a literal equation}),\\
 \llbracket Q_yB(y)\rrbracket
   &=\llbracket B(0)\rrbracket\to
                  \forall y\,\llbracket B(y)\rrbracket
       &&(B\text{ compound}),\\
 \llbracket D_{x,y}B\rrbracket
   &=\exists x\,\llbracket Q_yB\rrbracket.
 \end{aligned}
 \label{eq:new-reading}
\end{equation}
This translation preserves all free numerical variables and commutes with
numerical substitution.  It is a finite syntactic translation, not a truth
predicate on numerical formula codes.

\begin{theorem}[Soundness]
\label{thm:new-soundness}
If $\Gamma\vdash_{\mathcal R}F$, then
$\llbracket\Gamma\rrbracket\vdash_{T_{\rm eq}}\llbracket F\rrbracket$.
\end{theorem}
\begin{proof}
Hyp, computations, equational transport, and DE have their ordinary meanings.
Atomic DI is universal introduction of progress, followed by existential
introduction at its witness.  Compound DI is conditional universal
introduction using numerical coverage, followed by existential introduction.

For atomic QE, retain
$H=\forall y(E(y)\to E(Sy))$ and $E(a)$.
Equation induction on $n$ gives $E(\mathsf A(n,a))$: the base is $E(a)$,
and the step uses $H$ at $\mathsf A(n,a)$ and IS.  The pure equational
proofs in Lemma~\ref{lem:new-addition} also hold in $T_{\rm eq}$, so
commutativity gives $E(\mathsf A(a,n))$.  For compound QE the minor is
$\llbracket B(0)\rrbracket$; its conditional-universal major directly
supplies the target instance.  Thus every rule translates soundly.
\end{proof}

\subsection{A precise comparison with the existing native calculus}
\label{subsec:new-native-comparison}

Let $\mathcal N$ denote exactly the current native Eq/$Q$/$D$ calculus
\texttt{NPrf}.  Its $Q_yB$ always means $B(0)\to\forall yB(y)$; its direct
progress rule introduces $Q_yE$ only for equations.  Its terms and numerical
equality are identical to those of $\mathcal R$.

For $E\equiv(s=t)$ define
\begin{equation}
 \begin{aligned}
 C(v,a,b)&:=I(v,a;u.b),\\
 G_vE&:=\bigl(C(v,0,s)=C(v,0,t)\bigr),\\
 L(E)&:=D_{r,v}G_vE,\\
 P(F)&:=\begin{cases}L(F),&F\text{ an equation},\\F,&F\text{ compound}.
              \end{cases}
 \end{aligned}
 \label{eq:new-padding}
\end{equation}
All newly displayed binders are fresh; the update in $C$ ignores $u$, and
$r$ is unused in $L(E)$.

\begin{lemma}[Padding]
\label{lem:new-padding}
In $\mathcal R$, $L(E)$ and $E$ are interderivable under arbitrary contexts.
Consequently $P(F)$ is compound and interderivable with $F$.
\end{lemma}
\begin{proof}
Computations and equality prove $G_0E$ and
$G_{Sz}E\dashv\vdash E$.  From $E$, introduce $L(E)$ by atomic DI, ignoring
its local induction assumption and using $E$ for the successor target.
Conversely, DE on $L(E)$ gives $Q_vG_vE$.  Supply the theorem $G_0E$ and
apply QE with source $0$, initializer $1$.  I0 gives $G_1E$, which reduces
to $E$.  This proof uses no ordinary logical macro.
\end{proof}

Define a translation $T:\mathcal N\to\mathcal R$ by
\begin{equation}
 \begin{aligned}
 T(s=t)&=(s=t),\\
 T(Q_yB)&=Q_yP(TB),\\
 T(D_{x,y}B)&=D_{x,y}P(TB).
 \end{aligned}
 \label{eq:new-T}
\end{equation}
It fixes terms, commutes with numerical substitution, and places a compound
matrix immediately below every translated binder.

\begin{theorem}[Translation into the new calculus]
\label{thm:new-T}
Every $\mathcal N$ derivation translates effectively:
\[
 \Gamma\vdash_{\mathcal N}A
 \quad\Longrightarrow\quad
 T\Gamma\vdash_{\mathcal R}TA.
\]
\end{theorem}
\begin{proof}
Induct on the old derivation.  Hyp, computations, and equational transport
translate literally.  Old QE becomes compound new QE: pad the minor,
specialize at the old target using source $0$, remove the I0 redex, and
unpad.  In DE the opened new assumption is exactly
$Q_yP(TB)(b,y)$, the translation of the old opened $Q$ assumption.
For case DI, unpad the local base, apply the translated premise, and pad
the successor conclusion; the new matrix is compound.

For old direct Q-progress at $E$, the translated target is $Q_yL(E(y))$.
Introduce it by compound QI.  Its local base $L(E(0))$ supplies $E(0)$.
Apply the recovered contextual equation induction to the translated step,
then pad $E(Sz)$ to $L(E(Sz))$.  This discharges the only nonliteral rule
translation.  Each construction is finite and preserves eigenvariable
conditions after fresh renaming.
\end{proof}

\subsection{Logical macros and the order of the completeness proof}
\label{subsec:new-image-logic}

The old ordinary logical macros will be denoted
$\bot_{\mathcal N}$, $\to_{\mathcal N}$, $\wedge_{\mathcal N}$,
$\vee_{\mathcal N}$, $\exists_{\mathcal N}$, and $\forall_{\mathcal N}$.
For specificity they are the macros of the pinned
\texttt{NativeLogic.lean}:\cite{native}
\begin{equation}
 \begin{aligned}
 \bot_{\mathcal N}&=(0=1),\\
 A\to_{\mathcal N}B&=Q_y\operatorname{sel}(y,A,B),\\
 A\wedge_{\mathcal N}B
   &=Q_y\operatorname{gate}
           (y,\operatorname{sel}(\operatorname{par}(y),A,B)),\\
 \exists_{\mathcal N}x\,A(x)&=D_{x,y}\operatorname{gate}(y,A(x)),\\
 A\vee_{\mathcal N}B&=\exists_{\mathcal N}v\,
                                   \operatorname{sel}(v,A,B),\\
 \forall_{\mathcal N}x\,A(x)&=A(0)\wedge_{\mathcal N}Q_xA(x).
 \end{aligned}
 \label{eq:new-old-macros}
\end{equation}
Here $\operatorname{par}(y)=I(y,0;u.C(u,1,0))$.
The old gate acts on equations as $G_v$ in \eqref{eq:new-padding}.
The exact finite recursive definitions of the gate and selector, including
their behavior under $Q$ and $D$, are stated in Appendix~\ref{app:old-macros}.
The selector has zero endpoint $A$ and successor endpoint $B$.
All these operations concern the \emph{old} uniform conditional calculus.
The pinned repository supplies their capture-free definitions and derived
rules.

\begin{lemma}[Logic on image formulas, with arbitrary side context]
\label{lem:new-image-logic}
For old formulas $A,B$, the formulas
$T(A\to_{\mathcal N}B)$, $T(A\wedge_{\mathcal N}B)$, and the analogous
translated macros have their ordinary natural-deduction rules on principal
formulas $TA,TB$.  These derived rules permit an arbitrary retained
$\mathcal R$ context and arbitrary supplied $\mathcal R$ premise proofs.
Image existential elimination also permits an arbitrary new conclusion,
not necessarily an image formula.
\end{lemma}
\begin{proof}
Translate the finite \emph{derived-rule constructions}, rather than merely
closed old theorems.  Their nonlogical steps are replaced by the local
simulations in Theorem~\ref{thm:new-T}.  Their open premise leaves are filled
with the supplied new proofs; assumption cut and weakening
(Lemma~\ref{lem:new-structural}) retain any additional new context.  The
construction never examines the syntax of passive side assumptions.
Thus this is a uniform rule simulation, not an inference from
closed-theorem preservation alone.

For the important existential exception, the old elimination construction
opens its defining $D$ by DE, obtains the witness instance from the opened
gated closure, and cuts this instance into the supplied minor derivation.
It neither analyzes nor introduces the minor's conclusion.  The translated
construction consequently works for every new conclusion $F$, using the
same freshness condition.  This is precisely the structure of the pinned
\texttt{NPrf.exE} proof.
\end{proof}

We do \emph{not} yet assert that the old selectors or logical macros can be
applied unchanged to arbitrary mixed formulas.  Such an assertion would
incorrectly ignore the root distinction in $Q$.  The next normalization
argument first proves that every mixed formula is equivalent to an image.

\subsection{Both formula composites for the entire language}
\label{subsec:new-composites}

Define $J:\mathcal R\to\mathcal N$, fixing all terms and equations, by
\begin{equation}
 \begin{aligned}
 J(Q_yE)&=\forall_{\mathcal N}y\,
                (E(y)\to_{\mathcal N}E(Sy)),
                    &&E\text{ an equation},\\
 J(Q_yB)&=Q_yJB,     &&B\text{ compound},\\
 J(D_{x,y}E)&=\exists_{\mathcal N}x\,
             \forall_{\mathcal N}y\,
                (E(x,y)\to_{\mathcal N}E(x,Sy)),
                    &&E\text{ an equation},\\
 J(D_{x,y}B)&=D_{x,y}JB,&&B\text{ compound}.
 \end{aligned}
 \label{eq:new-J}
\end{equation}
This is a fixed finite syntax transformation.  It maps compound formulas to
compound formulas, preserves free variables, and commutes with numerical
substitution.  Its ordinary reading in $T_{\rm eq}$ is provably equivalent to
\eqref{eq:new-reading}, by the established meanings of the old logical macros.

\begin{lemma}[Atomic closure normalization]
\label{lem:new-atomic-normalization}
For every equation $E(x,y)$, under arbitrary retained contexts,
\begin{equation}
 \begin{aligned}
 Q_yE(x,y)&\dashv\vdash_{\mathcal R}
       T\bigl(\forall_{\mathcal N}y
           (E(x,y)\to_{\mathcal N}E(x,Sy))\bigr),\\
 D_{x,y}E(x,y)&\dashv\vdash_{\mathcal R}
       T\bigl(\exists_{\mathcal N}x\,
             \forall_{\mathcal N}y
           (E(x,y)\to_{\mathcal N}E(x,Sy))\bigr).
 \end{aligned}
 \label{eq:new-atomic-normalization}
\end{equation}
\end{lemma}
\begin{proof}
Write $H(x)$ for the image formula on the first line's right.
From $Q_yE(x,y)$ and the temporary assumption $E(x,z)$, derived one-step
elimination gives $E(x,Sz)$.  The image implication and universal
introductions of Lemma~\ref{lem:new-image-logic} give $H(x)$.  Their
principal component formulas are just equations, which are already fixed
by $T$; no full mixed-language logical rule has been assumed.
Conversely, under $H(x)$ and $E(x,z)$, image universal specialization and
implication elimination give $E(x,Sz)$; atomic QI gives $Q_yE(x,y)$.

For $D$, eliminate the left-hand binder at fresh $b$, normalize its local
$Q_yE(b,y)$ to $H(b)$ by the first part, introduce the image existential
at $b$, and close DE.  In the reverse direction, image existential
elimination supplies a fresh $b$ and $H(b)$.  Under $E(b,z)$ its image
universal and implication rules give $E(b,Sz)$, so atomic DI at witness
$b$ gives $D_{x,y}E$.  Close image existential elimination.  Its conclusion
is allowed to be this arbitrary new formula by the last assertion of
Lemma~\ref{lem:new-image-logic}.
\end{proof}

\begin{theorem}[Exact native-language round trips]
\label{thm:new-roundtrips}
For every old formula $A$ and every new formula $F$,
\begin{equation}
 A\dashv\vdash_{\mathcal N}JT(A),\qquad
 F\dashv\vdash_{\mathcal R}TJ(F).
 \label{eq:new-roundtrips}
\end{equation}
These equivalences are uniform under arbitrary additional retained contexts.
\end{theorem}
\begin{proof}
For the old-side composite, note first that $JL(E)$ is old-equivalent to
$E$.  It asserts existence of an unused witness and progressivity of the
guarded equation $G_vE$.  From $E$, every positive guard instance follows,
so the progress implications hold.  Conversely $G_0E$ is a theorem and
progress at $0$ yields $G_1E$, hence $E$.  This uses ordinary logical
macros in $\mathcal N$, whose correctness is already established.
Now induct on the old formula.  Every binder created by $T$ has a compound
matrix, so $J$ gives its old conditional interpretation.  The preceding
padding equivalence and the induction hypothesis apply under the old
$Q$ and $D$ congruence rules.

For the new-side composite, induct on $F$.  Equations are unchanged.
Atomic $Q$ and $D$ cases are exactly
Lemma~\ref{lem:new-atomic-normalization}.  Suppose the matrix $B$ is
compound.  The induction hypothesis gives $B\dashv\vdash TJB$.
Both $B$ and $TJB$ are compound, since $J$ and $T$ preserve compound roots.
Homogeneous congruence (Lemma~\ref{lem:new-transport}) therefore gives
\[
 Q B\dashv\vdash Q(TJB),\qquad
 D B\dashv\vdash D(TJB).
\]
Because $TJB$ is compound, no padding is inserted on these right sides;
they are exactly the required $TJ$ composites.  There is no appeal to
congruence across unlike root classes.  The induction proves normalization
for every formula before ordinary logic is extended to arbitrary new
formulas, avoiding a circular completeness argument.
\end{proof}

\paragraph{All ordinary logical operations in the new language.}
Now, and only now, define for arbitrary new formulas
\begin{equation}
 \begin{aligned}
 \bot_{\mathcal R}&=(0=1),\\
 F\mathbin{\diamond_{\mathcal R}}G
   &:=T\bigl(JF\mathbin{\diamond_{\mathcal N}}JG\bigr)
       &&(\diamond\in\{\to,\wedge,\vee\}),\\
 \exists_{\mathcal R}x\,F&:=T(\exists_{\mathcal N}x\,JF),\\
 \forall_{\mathcal R}x\,F&:=T(\forall_{\mathcal N}x\,JF).
 \end{aligned}
 \label{eq:new-full-macros}
\end{equation}
The round trips convert arbitrary new principal formulas to their images;
Lemma~\ref{lem:new-image-logic} then supplies the ordinary rules, and the
round trips convert their conclusions back.  For disjunction or existential
elimination with a new conclusion, normalize that conclusion first if
necessary.  All eigenvariable conditions are preserved because $J,T$
preserve free numerical variables.

This gives intuitionistic natural deduction for the entire new language,
including explosion, ordinary numerical coverage, and the quantified
rules.  It adds no primitive logical axioms.  In particular, the system does
not use excluded middle for arbitrary quantified formulas.

\subsection{Bi-interpretability, consistency, and the inherited arithmetic}
\label{subsec:new-final-bridge}

\begin{theorem}[Full consequence equivalence and bi-interpretability]
\label{thm:new-biinterpretability}
The translations $T,J$ preserve derivations in both directions and have the
round trips \eqref{eq:new-roundtrips}.  The ordinary realization of
$\mathcal R$ is equiconsistent and bi-interpretable with full first-order
iPRA.  Both numerical domains are whole, numerical equality is actual
equality, and the numerical composite isomorphisms have graph $x=y$.
\end{theorem}
\begin{proof}
Preservation by $T$ is Theorem~\ref{thm:new-T}.  For $J$, translate a new
derivation into $T_{\rm eq}$ using Theorem~\ref{thm:new-soundness}; its
formulas are provably equivalent to the ordinary readings of their $J$
translations.  The pinned
native completeness and native formula round-trip theorems translate this
ordinary derivation back into $\mathcal N$ and identify each resulting
formula with its $J$ translation.  Equivalently, one may translate each new
rule directly using the old ordinary macros and equation induction; the
only added arithmetic calculation is the QE soundness proof above.
The two exact formula composites are Theorem~\ref{thm:new-roundtrips}.

The remaining scalar-to-iPRA arithmetic is inherited, not reproved by the
new elimination rule.  The repository's chain constructs a scalar
predecessor, a term reflecting actual equality, quantifier-free induction,
pairing, and a compiler for each fixed primitive-recursive description.
The reverse interpretation expresses a fixed scalar iterator by primitive
recursion.  Equation induction proves the two numerical compilation round
trips.  These are the contents of the pinned \texttt{biInterp},
\texttt{native\_iff\_pra}, \texttt{pra\_iff\_native},
\texttt{native\_roundTrip}, and \texttt{pra\_native\_roundTrip} results.\cite{native,bridge}

Composing that existing bi-interpretation with $T,J$ proves the assertion.
Our two translations leave every scalar term and the numerical domain
unchanged, so they add neither a cut nor a quotient to the existing
identity graphs.  The designated contradiction $0=1$ is fixed.  Effective
finite proof translations therefore give equiconsistency.  This is an
external equivalence theorem, not an internal assertion of iPRA's own
consistency.
\end{proof}

\paragraph{What changes compared with the earlier progressive proposal.}
The repository already contains a presentation with a progressive binder
on every matrix, separate one-step and equation-iteration eliminations,
counter-aware primitive recursion, and an explicit Boolean-cases rule.\cite{progressive}
The progressive reading itself is therefore inherited motivation.  The
present contribution is the root-dependent complete reading with the scalar
iterator unchanged, the single reversed-orbit specialization
\eqref{eq:new-qe}, the removal of a separate case principle, and both
round trips for the entire Eq/$Q$/$D$ formula language.  The mixed reading
without the stronger elimination would lose information at atomic closures;
its soundness alone would not establish the theorem above.

\paragraph{Verification status.}
The new metatheorems in this section are written proofs.  They are not claimed
to have been checked in Lean or Rocq.  The existing Lean theorem names identify
the precise inherited bridge.  Executable certificates and checker tests
support concrete implementations; they do not replace the universal proof
or establish a global lower bound on presentation size.

% END reference/progressive/new_theory.tex


% BEGIN reference/progressive/old_macros.tex
\section{Explicit logical macros in the reference language}
\label{app:old-macros}

This appendix fixes the finite macros used by the interpretations. All
operators in this appendix belong to the old conditional language
$\mathcal N$. They must be expanded there before applying the translation
$T$. Applying these definitions unchanged to arbitrary new mixed
closures would be incorrect. These definitions reproduce the interface
of \texttt{NativeGate.lean} and \texttt{NativeLogic.lean}
\cite{native}.

Write $C(v,a,b)=I(v,a;u.b)$, where $u$ is unused, and $1=S0$.
Bindings below are freshened so that no displayed passive parameter is
captured. Define an equation--formula selector $U_v(s,t;B)$ by
\[
\begin{aligned}
U_v(s,t;r=w)&=(C(v,s,r)=C(v,t,w)),\\
U_v(s,t;Q_y B)&=Q_y U_v(C(y,0,s),C(y,0,t);B),\\
U_v(s,t;D_{x,y} B)&=D_{x,y}U_v(C(y,0,s),C(y,0,t);B).
\end{aligned}
\]
Its two endpoints are the equation $s=t$ and the formula $B$.
Put $G_v A=U_v(0,0;A)$. Then $G_0 A$ is provable and
$G_{Sz}A\eqv A$. A selector between arbitrary formulas is
\[
\begin{aligned}
M_v(s=t,B)&=U_v(s,t;B),\\
M_v(Q_y A,B)&=Q_y M_v(A,G_yB),\\
M_v(D_{x,y} A,B)&=D_{x,y}M_v(A,G_yB).
\end{aligned}
\]
This recursion terminates on the depth of its first argument; the
separate recursion defining $U$ terminates on its formula argument.
For every $A,B$, the old calculus derives
$M_0(A,B)\eqv A$ and $M_{Sz}(A,B)\eqv B$.

Let $p(y)=I(y,0;u.C(u,1,0))$. The logical operations are
\[
\begin{aligned}
\bot_*&=(0=1),\\
A\Rightarrow_* B&=Q_y M_y(A,B),\\
A\wedge_* B&=Q_y G_y M_{p(y)}(A,B),\\
\exists_* x\,A(x)&=D_{x,y}G_y A(x),\\
A\vee_*B&=\exists_*v\,M_v(A,B),\\
\forall_*x\,A(x)&=A(0)\wedge_* Q_x A(x).
\end{aligned}
\]
All new bound variables are fresh. These are definitions by finite
syntax, not additional primitive logical axioms or rules.

\paragraph{Why they have the logical rules.}
Implication introduction uses the zero selector as its local hypothesis
and the successor selector as its conclusion; elimination specializes
at $1$. Existential introduction chooses the witness and uses the
successor gate. Existential elimination uses $D$ elimination, the
zero-gate theorem, and specialization at $1$; its final conclusion can
be any formula, independently of the existential matrix.

For conjunction, the gate supplies the base theorem and the computations
$p(1)=1$, $p(2)=0$ recover the two conjuncts. Numerical cases establish
the selector at an arbitrary control from proofs of both alternatives.
Disjunction uses witnesses $0$ and $1$; its elimination uses the derived
case rule on an implication to the desired conclusion. The universal
macro has its zero instance as one conjunct and every instance as the
other conjunct's closure. Finally, $0=1$ and the selector computations
give $s=C(0,s,t)=C(1,s,t)=t$ for arbitrary terms; structural induction
on formulas gives falsity elimination.

The full proofs and their substitution lemmas are preserved in the
pinned source files included with the package. The new theory's
normalization proof uses the translated versions of these derived-rule
constructions on $T$-image principal formulas before it defines the
ordinary logical operations on arbitrary new formulas.

% END reference/progressive/old_macros.tex


% BEGIN reference/progressive/arithmetic_bridge.tex
\section{The arithmetic interface and the inherited bridge}
\label{app:arithmetic-bridge}

The reference theory $\Teq$ is intuitionistic first-order equality logic
over scalar iterator terms, with the two computation equations,
numerical coverage
\[
\forall x\,(x=0\vee\exists y\,x=Sy),
\]
separation $(0=S0)\to\bot$, and induction for literal equations under
side assumptions independent of the induction variable. Its equality
is ordinary equality. The reference iPRA has a symbol for each fixed
primitive recursive definition, its defining equations, successor
separation and injectivity, and quantifier-free induction over ordinary
intuitionistic first-order logic.

The pinned development establishes $\mathcal N\simeq\Teq\simeq\ipra$
\cite{bridge}. The arithmetic part is retained unchanged here. Its
substantive stages are:
\begin{enumerate}
\item Construct a predecessor in the scalar iterator language. Since the
update binds only its previous state, this is a real construction; the
counter variable is not silently made available to the update.
\item Derive equality reflection into a term and Boolean operations.
Replace a quantifier-free formula by its characteristic equation, which
transfers literal-equation induction to quantifier-free induction.
\item Construct pairing and projections, and compile a counter-aware
primitive recursion by iterating a state that carries both the counter
and the accumulated value.
\item Translate each fixed primitive recursive definition and each
scalar iterator term in the opposite direction, proving their defining
equations and both composites by equation induction.
\end{enumerate}

The archive contains the corresponding complete written sources
\path{scalar_predecessor.tex}, \path{equality_reflection.tex}, and
\path{primitive_recursion.tex} in the pinned reference subtree, together
with the ordinary bridge and the native logical basis. These files are
included to expose the mathematical dependency; the new report does
not claim to have discovered these arithmetic constructions.

There is no quantified code for an arbitrary primitive recursive
program and no total universal primitive recursive evaluator in the
theory. A compiler acts externally on one finite definition at a time.
This distinction matters: adjoining a purported total universal
evaluator would change the foundational argument.

Both new translations between $\mathcal R$ and $\mathcal N$ fix every
numerical term and literal equation. Composing them with the reference
arithmetic translations therefore retains whole numerical domains and
actual equality. The graph of each numerical interpretation composite
is provably $x=y$. Formula composites are proved for all formulas, so
the conclusion is bi-interpretability, not merely a correspondence of
closed equations or an equiconsistency argument.

The precise existing Lean endpoint names are
\begin{quote}\small\ttfamily
native\_iff\_pra,\quad pra\_iff\_native,\\
native\_roundTrip,\quad pra\_native\_roundTrip,\\
native\_inconsistent\_iff.
\end{quote}
Their source was retrieved and
authenticated at the pinned commit. Lean/Lake was not available in this
run, so these existing files were inspected rather than rebuilt, and no
new Lean theorem about $\mathcal R$ is claimed.

% END reference/progressive/arithmetic_bridge.tex

\begin{thebibliography}{99}
\bibitem{native} ProveIt, \emph{Native scalar syntax, gates, and logic},
\texttt{Native.lean}, \texttt{NativeGate.lean}, and \texttt{NativeLogic.lean},
at the pinned commit.
\bibitem{bridge} ProveIt, \emph{NativeBiInterpretation.lean} and
\emph{BiInterpretation.lean}, at the pinned commit. The endpoint sources
are retained with the archive.
\bibitem{generated} ProveIt research series, \emph{Complementary introductions
with one local assumption}, retained in
\path{reference/foundations/inherited_choice_basis.tex}.
\bibitem{erasure} ProveIt research series, scalar predecessor construction,
retained in \path{reference/foundations/scalar_predecessor.tex}.
\bibitem{forward} ProveIt research series, equality reflection,
retained in \path{reference/foundations/equality_reflection.tex}.
\bibitem{progress} ProveIt research series, full primitive recursion,
retained in \path{reference/foundations/primitive_recursion.tex}.
\bibitem{source} ProveIt, \emph{source\_dispatch},
\path{proofs/zero_source_admissibility.tex}, pinned commit.
\bibitem{witness} ProveIt, \emph{NativeOneIntro.lean}, witness-sensitive
single-introduction equivalence theorem, pinned repository.
\bibitem{progressive} ProveIt, \emph{Research/research/theory/progressive\_binder.tex},
the earlier progressive-binder development cited by the inherited report.
\bibitem{uniform} ProveIt incoming research, \emph{A uniform introduction
rule for intuitionistic primitive recursive arithmetic}, archive
\texttt{ipra\_uniform\_progress\_2026-11-11.zip}.
\bibitem{repo} V. Reshetnikov and ProveIt contributors, \emph{ConciseIPRA}.
\url{https://github.com/VladimirReshetnikov/ProveIt/tree/76697262071ada91f29c14d9cf3801a64dacb994/Logic/PrimitiveRecursiveArithmetic/ConciseIPRA}.
\end{thebibliography}
\end{document}
